\documentclass[11pt]{article}
\usepackage{../common/ai204}
\usepackage{animate}
\renewcommand{\doctitle}{Laboratory Manual}
\begin{document}
\aicover{Laboratory Manual}{10 experiments \textbullet\ Python 3 \textbullet\ executed Jupyter notebooks with outputs \textbullet\ animated algorithm visualisations}{For DSY -- AI \& ML students \textbullet\ Integrated Theory + Practical (2-0-2)}
\section*{Certificate}
\begin{tcolorbox}[enhanced,colback=ivory,colframe=gold,boxrule=1.2pt,arc=3mm,left=8mm,right=8mm,top=6mm,bottom=6mm]
\large This is to certify that Mr./Ms.\ \rule{7cm}{0.4pt}\ Roll No.\ \rule{2.5cm}{0.4pt}\ of class \rule{3cm}{0.4pt}\ has satisfactorily completed the experiments in \textbf{26UAMPC204 Artificial Intelligence Laboratory} during the academic year 2026--27.\\[3cm]
\rule{4.5cm}{0.4pt}\hfill\rule{4.5cm}{0.4pt}\hfill\rule{4.5cm}{0.4pt}\\ Course Faculty\hfill Head of Department\hfill External Examiner
\end{tcolorbox}
\section*{Laboratory Instructions}
\begin{itemize}\item Install Python 3.10+ and run \texttt{pip install -r requirements.txt} from the repository root.
\item Open each notebook in Jupyter or VS Code, run all cells, and compare your output with the output printed in this manual.
\item Write the aim, algorithm, program, output and result in your journal; attach the printed figures.
\item Each experiment carries 10 marks: performance 4, journal 3, viva 3.
\item Animations play inside this PDF in Adobe Acrobat Reader; GIF versions are in \texttt{assets/animations}.\end{itemize}
\section*{Index of Experiments}
\begin{center}\begin{tabular}{c p{8.6cm} c c p{1.6cm} p{1.6cm}}
\rowcolor{navy}\textcolor{gold}{\textbf{No.}}&\textcolor{gold}{\textbf{Title}}&\textcolor{gold}{\textbf{CO}}&\textcolor{gold}{\textbf{Page}}&\textcolor{gold}{\textbf{Date}}&\textcolor{gold}{\textbf{Sign}}\\
1&Uninformed Search --- BFS and DFS&CO204.1&\pageref{exp1}&&\\
\rowcolor{ivory}2&A* Search with Admissible Heuristics&CO204.2&\pageref{exp2}&&\\
3&Hill Climbing and Greedy Best-First Search&CO204.2&\pageref{exp3}&&\\
\rowcolor{ivory}4&Game Playing --- Minimax with Alpha--Beta Pruning&CO204.2&\pageref{exp4}&&\\
5&Constraint Satisfaction --- N-Queens, Map Colouring and Sudoku&CO204.2&\pageref{exp5}&&\\
\rowcolor{ivory}6&Facts, Rules and Logical Inference&CO204.3&\pageref{exp6}&&\\
7&Unification and Resolution in First-Order Logic&CO204.3&\pageref{exp7}&&\\
\rowcolor{ivory}8&Bayes' Theorem and a Bayesian Network for Diagnosis&CO204.4&\pageref{exp8}&&\\
9&Decision-Tree Classifier and Rule Extraction&CO204.5&\pageref{exp9}&&\\
\rowcolor{ivory}10&Rule-Based Expert System and a Basic Chatbot (Mini-Project)&CO204.5&\pageref{exp10}&&\\
\end{tabular}\end{center}
\section*{Course Outcomes}
\begin{tabularx}{\linewidth}{>{\bfseries\color{navy}}l X}\toprule
CO204.1&Explain the foundations of AI and intelligent agents and formulate problems as state-space search.\\
CO204.2&Apply uninformed, informed, local and adversarial search, and constraint satisfaction.\\
CO204.3&Represent knowledge and perform inference in propositional and first-order logic.\\
CO204.4&Apply classical planning and reasoning under uncertainty (probability, Bayesian networks, fuzzy logic).\\
CO204.5&Illustrate learning, expert systems and applications of AI with awareness of ethical issues.\\\bottomrule
\end{tabularx}
\newpage
\section*{Experiment 1: Uninformed Search --- BFS and DFS}\label{exp1}
\noindent\co{CO204.1}\ \btl{L3 Apply}\ \btl{L4 Analyse}\hfill\textbf{\color{navy}Date:}\ \rule{3cm}{0.4pt}\quad\textbf{\color{navy}Sign:}\ \rule{3cm}{0.4pt}\par\medskip
\begin{keypoint}[Aim]Implement Breadth-First Search and Depth-First Search to solve the Water-Jug problem and the 8-puzzle, and compare nodes expanded, path length and memory.\end{keypoint}
\subsection*{Theory}
\textbf{State space.} A problem is defined by an initial state, actions, a transition model, a goal test and path cost.

\textbf{BFS} expands the shallowest node first using a FIFO queue --- complete and optimal for unit costs, but needs $O(b^d)$ memory.

\textbf{DFS} expands the deepest node first using a LIFO stack --- needs only $O(bm)$ memory but is neither complete (in infinite spaces) nor optimal.

\textbf{Water-Jug (4 L, 3 L):} state $(x,y)$; actions fill, empty and pour; goal $x=2$.
\subsection*{Algorithm}
\begin{algo}[Procedure]\begin{enumerate}
\item Put the start node in the frontier (queue for BFS / stack for DFS) and mark it reached.
\item Pop a node; if it satisfies the goal test, reconstruct the path through parent links.
\item Otherwise generate successors; push every successor not yet reached and record its parent.
\item Repeat until the goal is found or the frontier is empty.
\end{enumerate}\end{algo}
\subsection*{Program (Python 3)}
{\small\color{slate}Notebook: \texttt{06\_Notebooks/Exp01\_BFS\_DFS\_WaterJug\_8Puzzle.ipynb}\\ Script: \texttt{06\_Notebooks/python/Exp01\_BFS\_DFS\_WaterJug\_8Puzzle.py}}
\begin{lstlisting}[style=py]
from collections import deque

CAP = (4, 3)            # jug capacities
START, GOAL_X = (0, 0), 2

def successors(s):
    x, y = s; A, B = CAP
    moves = {
        "Fill 4L":  (A, y), "Fill 3L": (x, B),
        "Empty 4L": (0, y), "Empty 3L": (x, 0),
        "Pour 4->3": (x - min(x, B - y), y + min(x, B - y)),
        "Pour 3->4": (x + min(y, A - x), y - min(y, A - x)),
    }
    return [(a, n) for a, n in moves.items() if n != s]

def is_goal(s): return s[0] == GOAL_X
print("Successors of (0,0):", successors((0, 0)))
\end{lstlisting}
\noindent{\small\bfseries\color{gold!70!black}Output}
\begin{lstlisting}[style=out]
Successors of (0,0): [('Fill 4L', (4, 0)), ('Fill 3L', (0, 3))]
\end{lstlisting}
\begin{lstlisting}[style=py]
def search(start, succ, goal, strategy="BFS", limit=10**6):
    frontier = deque([start]); parent = {start: (None, None)}
    order, max_frontier = [], 1
    while frontier:
        s = frontier.popleft() if strategy == "BFS" else frontier.pop()
        order.append(s)
        if goal(s):
            path, acts = [s], []
            while parent[path[-1]][0] is not None:
                p, a = parent[path[-1]]; acts.append(a); path.append(p)
            return path[::-1], acts[::-1], order, max_frontier
        for a, n in succ(s):
            if n not in parent:
                parent[n] = (s, a); frontier.append(n)
        max_frontier = max(max_frontier, len(frontier))
        if len(order) > limit: break
    return None, None, order, max_frontier

for strat in ("BFS", "DFS"):
    path, acts, order, mf = search(START, successors, is_goal, strat)
    print(f"\n{strat}: path length = {len(path)-1}, nodes expanded = {len(order)}, max frontier = {mf}")
    for s, a in zip(path, ["Start"] + acts):
        print(f"   {a:<9} -> {s}")
\end{lstlisting}
\noindent{\small\bfseries\color{gold!70!black}Output}
\begin{lstlisting}[style=out]
BFS: path length = 6, nodes expanded = 13, max frontier = 3
   Start     -> (0, 0)
   Fill 4L   -> (4, 0)
   Pour 4->3  -> (1, 3)
   Empty 3L  -> (1, 0)
   Pour 4->3  -> (0, 1)
   Fill 4L   -> (4, 1)
   Pour 4->3  -> (2, 3)

DFS: path length = 6, nodes expanded = 7, max frontier = 3
   Start     -> (0, 0)
   Fill 3L   -> (0, 3)
   Pour 3->4  -> (3, 0)
   Fill 3L   -> (3, 3)
   Pour 3->4  -> (4, 2)
   Empty 4L  -> (0, 2)
   Pour 3->4  -> (2, 0)
\end{lstlisting}
\begin{lstlisting}[style=py]
import networkx as nx
G = nx.DiGraph()
for x in range(5):
    for y in range(4):
        for a, n in successors((x, y)): G.add_edge((x, y), n)
reach = nx.descendants(G, START) | {START}
H = G.subgraph(reach)
pos = {n: (n[0] * 1.4, n[1] * 1.2) for n in H.nodes}
_, _, bfs_order, _ = search(START, successors, is_goal, "BFS")

fig, ax = plt.subplots(figsize=(7.5, 5))
def frame(k):
    ax.clear(); seen = set(bfs_order[:k + 1])
    cols = [GOLD if n == bfs_order[k] else (TEAL if n in seen else MIST) for n in H.nodes]
    nx.draw_networkx_edges(H, pos, ax=ax, edge_color="#C8D0DC", arrows=True, arrowsize=8, width=0.8)
    nx.draw_networkx_nodes(H, pos, ax=ax, node_color=cols, edgecolors=NAVY, node_size=900)
    nx.draw_networkx_labels(H, pos, ax=ax, labels={n: f"{n[0]},{n[1]}" for n in H.nodes}, font_size=9, font_family="serif")
    ax.set_title(f"BFS step {k+1}: expanding {bfs_order[k]}", color=NAVY); ax.axis("off")
anim = animation.FuncAnimation(fig, frame, frames=len(bfs_order))
gif = save_gif(anim, "exp01_bfs_waterjug.gif", fps=1.5)
display(Image(filename=gif))
\end{lstlisting}
\begin{lstlisting}[style=py]
GOAL8 = (1, 2, 3, 4, 5, 6, 7, 8, 0)
def succ8(s):
    i = s.index(0); r, c = divmod(i, 3); out = []
    for a, dr, dc in (("Up", -1, 0), ("Down", 1, 0), ("Left", 0, -1), ("Right", 0, 1)):
        nr, nc = r + dr, c + dc
        if 0 <= nr < 3 and 0 <= nc < 3:
            j = nr * 3 + nc; t = list(s); t[i], t[j] = t[j], t[i]; out.append((a, tuple(t)))
    return out
rng8 = np.random.default_rng(11); start8, prev = GOAL8, None
for _ in range(9):                     # scramble the goal with 9 random non-reversing moves
    opts = [n for a, n in succ8(start8) if n != prev]; prev, start8 = start8, opts[rng8.integers(len(opts))]
print("Start state:", start8)
p, a, o, mf = search(start8, succ8, lambda s: s == GOAL8, "BFS")
print("BFS  : moves =", a, "| expanded =", len(o), "| max frontier =", mf)
p2, a2, o2, mf2 = search(start8, succ8, lambda s: s == GOAL8, "DFS", limit=20000)
print("DFS  :", "gave up after" if p2 is None else f"path length {len(p2)-1},", f"{len(o2):,} expansions (DFS wanders deep -- not optimal)")

def show(s, ax, title):
    ax.imshow(np.zeros((3, 3)), cmap="Greys", vmin=0, vmax=1)
    for k, v in enumerate(s):
        r, c = divmod(k, 3)
        ax.add_patch(plt.Rectangle((c - .45, r - .45), .9, .9, color=NAVY if v else "white", ec=GOLD, lw=2))
        if v: ax.text(c, r, v, ha="center", va="center", color=GOLD, fontsize=20, weight="bold")
    ax.set_title(title, fontsize=11); ax.axis("off")
fig, axs = plt.subplots(1, len(p), figsize=(2.1 * len(p), 2.4))
for ax, s, t in zip(axs, p, ["Start"] + a): show(s, ax, t)
plt.tight_layout(); plt.savefig(f"{FIG}/exp01_8puzzle_path.png", dpi=130); plt.show()
\end{lstlisting}
\noindent{\small\bfseries\color{gold!70!black}Output}
\begin{lstlisting}[style=out]
Start state: (4, 1, 2, 7, 5, 0, 8, 6, 3)
BFS  : moves = ['Down', 'Left', 'Left', 'Up', 'Up', 'Right', 'Right', 'Down', 'Down'] | expanded = 390 | 
DFS  : path length 319, 326 expansions (DFS wanders deep -- not optimal)
\end{lstlisting}
\begin{lstlisting}[style=py]
import pandas as pd
rows = []
for strat in ("BFS", "DFS"):
    pth, _, od, mfr = search(START, successors, is_goal, strat)
    rows.append([strat, len(pth) - 1, len(od), mfr])
df = pd.DataFrame(rows, columns=["Strategy", "Path length", "Nodes expanded", "Max frontier"])
print(df.to_string(index=False))
ax = df.set_index("Strategy")[["Path length", "Nodes expanded", "Max frontier"]].plot.bar(color=[NAVY, GOLD, TEAL], figsize=(6, 3.5), rot=0)
ax.set_title("Water-Jug: BFS vs DFS"); plt.tight_layout(); plt.savefig(f"{FIG}/exp01_compare.png", dpi=130); plt.show()
\end{lstlisting}
\noindent{\small\bfseries\color{gold!70!black}Output}
\begin{lstlisting}[style=out]
Strategy  Path length  Nodes expanded  Max frontier
     BFS            6              13             3
     DFS            6               7             3
\end{lstlisting}
\subsection*{Output Figures and Animation}
\begin{center}
\animategraphics[controls,loop,autoplay,width=0.62\linewidth]{2}{frames/exp01_bfs_waterjug/f-}{0}{12}\\
{\small\color{slate}\textbf{Animation} (13 frames; plays in Adobe Acrobat Reader). GIF: \texttt{assets/animations/exp01\_bfs\_waterjug.gif}}
\end{center}
\begin{center}
\includegraphics[width=0.78\linewidth,height=0.32\textheight,keepaspectratio]{../assets/figures/exp01_8puzzle_path.png}
\end{center}
\begin{center}
\includegraphics[width=0.78\linewidth,height=0.32\textheight,keepaspectratio]{../assets/figures/exp01_compare.png}
\end{center}
\subsection*{Result}
\begin{outcome}[Result \& Conclusion]Both strategies solved the Water-Jug problem in 6 moves, but only BFS guarantees the shortest solution. On the scrambled 8-puzzle BFS found the optimal path while DFS wandered thousands of nodes deep without returning a short solution. BFS suits shallow solutions; DFS suits tight memory.\end{outcome}
\subsection*{Viva-Voce Questions}
\begin{enumerate}
\item Why is BFS optimal only when all step costs are equal?
\item What is the space complexity of DFS and why is it smaller than BFS?
\item How do you detect that an 8-puzzle state is unsolvable?
\item What is the difference between tree search and graph search?
\item Why does the Water-Jug state graph contain cycles?
\end{enumerate}
\newpage
\section*{Experiment 2: A* Search with Admissible Heuristics}\label{exp2}
\noindent\co{CO204.2}\ \btl{L3 Apply}\ \btl{L4 Analyse}\hfill\textbf{\color{navy}Date:}\ \rule{3cm}{0.4pt}\quad\textbf{\color{navy}Sign:}\ \rule{3cm}{0.4pt}\par\medskip
\begin{keypoint}[Aim]Implement A* search with admissible heuristics for the 8-puzzle (misplaced tiles, Manhattan distance) and for a grid shortest-route problem, and compare with uniform-cost search.\end{keypoint}
\subsection*{Theory}
A* expands the node with the smallest $f(n)=g(n)+h(n)$. If $h$ is \textbf{admissible} ($h\le h^*$) tree-A* is optimal; if $h$ is \textbf{consistent} graph-A* is optimal.

For the 8-puzzle, $h_1$ = misplaced tiles and $h_2$ = Manhattan distance; $h_2$ dominates $h_1$, so it expands fewer nodes. With $h=0$, A* reduces to uniform-cost search.
\subsection*{Algorithm}
\begin{algo}[Procedure]\begin{enumerate}
\item Insert start with g = 0 and f = h(start) into a priority queue.
\item Pop the node with the lowest f; if it is the goal, return the path.
\item For each successor compute g' = g + step cost; if g' improves the best known g, record parent and push with f = g' + h.
\item Repeat until goal popped or queue empty.
\end{enumerate}\end{algo}
\subsection*{Program (Python 3)}
{\small\color{slate}Notebook: \texttt{06\_Notebooks/Exp02\_AStar\_8Puzzle\_Route.ipynb}\\ Script: \texttt{06\_Notebooks/python/Exp02\_AStar\_8Puzzle\_Route.py}}
\begin{lstlisting}[style=py]
import heapq, itertools, time
GOAL8 = (1, 2, 3, 4, 5, 6, 7, 8, 0)
def succ8(s):
    i = s.index(0); r, c = divmod(i, 3)
    for a, dr, dc in (("Up", -1, 0), ("Down", 1, 0), ("Left", 0, -1), ("Right", 0, 1)):
        nr, nc = r + dr, c + dc
        if 0 <= nr < 3 and 0 <= nc < 3:
            j = nr * 3 + nc; t = list(s); t[i], t[j] = t[j], t[i]; yield a, tuple(t), 1

def h0(s): return 0
def h1(s): return sum(1 for k, v in enumerate(s) if v and v != GOAL8[k])
def h2(s):
    d = 0
    for k, v in enumerate(s):
        if v: gr, gc = divmod(v - 1, 3); r, c = divmod(k, 3); d += abs(r - gr) + abs(c - gc)
    return d

def astar(start, succ, goal, h):
    tie = itertools.count(); frontier = [(h(start), next(tie), start)]
    g = {start: 0}; parent = {start: None}; expanded = 0
    while frontier:
        f, _, s = heapq.heappop(frontier)
        if f > g[s] + h(s): continue
        expanded += 1
        if goal(s):
            path = [s]
            while parent[path[-1]]: path.append(parent[path[-1]][0])
            return path[::-1], g[s], expanded
        for a, n, c in succ(s):
            ng = g[s] + c
            if ng < g.get(n, float("inf")):
                g[n] = ng; parent[n] = (s, a); heapq.heappush(frontier, (ng + h(n), next(tie), n))
    return None, None, expanded

start = (7, 2, 4, 5, 0, 6, 8, 3, 1)   # Russell & Norvig Fig. 3.25 start state
print("h1(start) =", h1(start), "  h2(start) =", h2(start))
\end{lstlisting}
\noindent{\small\bfseries\color{gold!70!black}Output}
\begin{lstlisting}[style=out]
h1(start) = 6   h2(start) = 14
\end{lstlisting}
\begin{lstlisting}[style=py]
import pandas as pd
res = []
for name, h in (("UCS (h=0)", h0), ("A* h1 misplaced", h1), ("A* h2 Manhattan", h2)):
    t = time.time(); path, cost, exp = astar(start, succ8, lambda s: s == GOAL8, h)
    res.append([name, cost, exp, round(time.time() - t, 2)])
df = pd.DataFrame(res, columns=["Algorithm", "Solution cost", "Nodes expanded", "Time (s)"])
print(df.to_string(index=False))
fig, ax = plt.subplots(figsize=(6.5, 3.5))
ax.barh(df["Algorithm"], df["Nodes expanded"], color=[SLATE, GOLD, NAVY]); ax.set_xscale("log")
for i, v in enumerate(df["Nodes expanded"]): ax.text(v, i, f"  {v:,}", va="center", color=NAVY)
ax.set_title("Nodes expanded (log scale) -- better heuristic, less work"); ax.invert_yaxis()
plt.tight_layout(); plt.savefig(f"{FIG}/exp02_heuristics.png", dpi=130); plt.show()
\end{lstlisting}
\noindent{\small\bfseries\color{gold!70!black}Output}
\begin{lstlisting}[style=out]
Algorithm  Solution cost  Nodes expanded  Time (s)
      UCS (h=0)             20           55409      0.24
A* h1 misplaced             20            3667      0.02
A* h2 Manhattan             20             283      0.00
\end{lstlisting}
\begin{lstlisting}[style=py]
rng = np.random.default_rng(7)
N = 20; grid = (rng.random((N, N)) < 0.25).astype(int)
S, T = (0, 0), (N - 1, N - 1); grid[S] = grid[T] = 0
def gsucc(p):
    for dr, dc in ((1,0),(-1,0),(0,1),(0,-1)):
        r, c = p[0] + dr, p[1] + dc
        if 0 <= r < N and 0 <= c < N and not grid[r, c]: yield "m", (r, c), 1
def manh(p): return abs(p[0] - T[0]) + abs(p[1] - T[1])

def astar_trace(h):
    tie = itertools.count(); fr = [(h(S), next(tie), S)]; g = {S: 0}; par = {S: None}; closed = []
    while fr:
        f, _, s = heapq.heappop(fr)
        if s in closed: continue
        closed.append(s)
        if s == T: break
        for _, n, c in gsucc(s):
            if g[s] + c < g.get(n, 1e9): g[n] = g[s] + c; par[n] = s; heapq.heappush(fr, (g[n] + h(n), next(tie), n))
    path = [T]
    while par.get(path[-1]): path.append(par[path[-1]])
    return closed, path[::-1]
closed_a, path_a = astar_trace(manh)
closed_u, path_u = astar_trace(lambda p: 0)
print(f"A*  : expanded {len(closed_a)} cells, path length {len(path_a)-1}")
print(f"UCS : expanded {len(closed_u)} cells, path length {len(path_u)-1}")

fig, ax = plt.subplots(figsize=(5.5, 5.5))
step = max(1, len(closed_a) // 30)
frames = list(range(0, len(closed_a), step)) + [len(closed_a)] * 4
def draw(k):
    ax.clear(); img = np.ones((N, N, 3))
    img[grid == 1] = matplotlib.colors.to_rgb(NAVY)
    for (r, c) in closed_a[:k]: img[r, c] = matplotlib.colors.to_rgb(GOLDL)
    if k >= len(closed_a):
        for (r, c) in path_a: img[r, c] = matplotlib.colors.to_rgb(GOLD)
    img[S] = matplotlib.colors.to_rgb(LEAF); img[T] = matplotlib.colors.to_rgb(CRIMSON)
    ax.imshow(img); ax.set_xticks([]); ax.set_yticks([])
    ax.set_title(f"A* (Manhattan): {min(k, len(closed_a))} cells expanded")
gif = save_gif(animation.FuncAnimation(fig, draw, frames=frames), "exp02_astar_grid.gif", fps=4)
display(Image(filename=gif))
\end{lstlisting}
\noindent{\small\bfseries\color{gold!70!black}Output}
\begin{lstlisting}[style=out]
A*  : expanded 134 cells, path length 38
UCS : expanded 287 cells, path length 38
\end{lstlisting}
\subsection*{Output Figures and Animation}
\begin{center}
\animategraphics[controls,loop,autoplay,width=0.62\linewidth]{2}{frames/exp02_astar_grid/f-}{0}{14}\\
{\small\color{slate}\textbf{Animation} (15 frames; plays in Adobe Acrobat Reader). GIF: \texttt{assets/animations/exp02\_astar\_grid.gif}}
\end{center}
\begin{center}
\includegraphics[width=0.78\linewidth,height=0.32\textheight,keepaspectratio]{../assets/figures/exp02_heuristics.png}
\end{center}
\subsection*{Result}
\begin{outcome}[Result \& Conclusion]All three searches return the optimal 20-move 8-puzzle solution, but A* with Manhattan distance expands orders of magnitude fewer nodes than UCS. On the grid, A* explores far fewer cells than UCS while returning a path of the same optimal length.\end{outcome}
\subsection*{Viva-Voce Questions}
\begin{enumerate}
\item Prove that the Manhattan distance heuristic is admissible.
\item What is meant by one heuristic dominating another?
\item When does A* become uniform-cost search?
\item What happens to A* if the heuristic overestimates?
\item Define consistency and explain why it implies admissibility.
\end{enumerate}
\newpage
\section*{Experiment 3: Hill Climbing and Greedy Best-First Search}\label{exp3}
\noindent\co{CO204.2}\ \btl{L3 Apply}\ \btl{L4 Analyse}\hfill\textbf{\color{navy}Date:}\ \rule{3cm}{0.4pt}\quad\textbf{\color{navy}Sign:}\ \rule{3cm}{0.4pt}\par\medskip
\begin{keypoint}[Aim]Implement steepest-ascent hill climbing (with random restarts) for the 8-queens problem and Greedy Best-First Search for a route problem; study local maxima and success rates.\end{keypoint}
\subsection*{Theory}
\textbf{Hill climbing} keeps one current state and moves to the best neighbour while it improves the objective. It can get stuck at \textbf{local maxima}, \textbf{plateaux} and \textbf{ridges}; \textbf{random restarts} make it complete with probability approaching 1.

\textbf{Greedy Best-First Search} expands the node with the lowest $h(n)$ --- fast, but not optimal.

For 8-queens, the objective counts attacking pairs $h$ (goal $h=0$).
\subsection*{Algorithm}
\begin{algo}[Procedure]\begin{enumerate}
\item Start from a random complete state.
\item Evaluate all neighbours (move one queen within its column).
\item If the best neighbour is not better than the current state, stop (local optimum).
\item Otherwise move to it and repeat; on failure, restart from a new random state.
\end{enumerate}\end{algo}
\subsection*{Program (Python 3)}
{\small\color{slate}Notebook: \texttt{06\_Notebooks/Exp03\_HillClimbing\_BestFirst.ipynb}\\ Script: \texttt{06\_Notebooks/python/Exp03\_HillClimbing\_BestFirst.py}}
\begin{lstlisting}[style=py]
rng = np.random.default_rng(42)
def attacks(s):
    n = len(s); return sum(1 for i in range(n) for j in range(i + 1, n) if s[i] == s[j] or abs(s[i] - s[j]) == j - i)
def hill_climb(n=8, max_steps=100):
    s = list(rng.integers(0, n, n)); trace = [s[:]]
    for _ in range(max_steps):
        cur = attacks(s); best, best_s = cur, None
        for col in range(n):
            for row in range(n):
                if row != s[col]:
                    t = s[:]; t[col] = row; v = attacks(t)
                    if v < best: best, best_s = v, t
        if best_s is None: break
        s = best_s; trace.append(s[:])
    return s, attacks(s), trace

wins, steps = 0, []
for trial in range(500):
    s, h, tr = hill_climb()
    wins += (h == 0); steps.append(len(tr) - 1)
print(f"Steepest-ascent success rate on 8-queens: {wins/500:.1%} (textbook ~= 14%)   mean steps = {np.mean(steps):.1f}")
\end{lstlisting}
\noindent{\small\bfseries\color{gold!70!black}Output}
\begin{lstlisting}[style=out]
Steepest-ascent success rate on 8-queens: 15.4% (textbook ~= 14%)   mean steps = 3.2
\end{lstlisting}
\begin{lstlisting}[style=py]
def restart_hc(n=8):
    restarts = 0
    while True:
        s, h, tr = hill_climb(n)
        if h == 0: return s, restarts, tr
        restarts += 1
sol, r, trace = restart_hc()
sol = [int(v) for v in sol]; trace = [[int(v) for v in t] for t in trace]
print("Solution:", sol, "found after", r, "random restarts;  attacking pairs =", attacks(sol))

def board(ax, s, title):
    n = len(s)
    ax.imshow([[(i + j) % 2 for j in range(n)] for i in range(n)], cmap=matplotlib.colors.ListedColormap([MIST, "#C8D0DC"]))
    for c, r in enumerate(s): ax.text(c, r, "Q", ha="center", va="center", fontsize=24, fontfamily="DejaVu Sans", color=NAVY if attacks(s) else LEAF)
    ax.set_title(title); ax.set_xticks([]); ax.set_yticks([])
fig, ax = plt.subplots(figsize=(4.5, 4.8))
gif = save_gif(animation.FuncAnimation(fig, lambda k: (ax.clear(), board(ax, trace[min(k, len(trace)-1)], f"Step {min(k, len(trace)-1)}: h = {attacks(trace[min(k, len(trace)-1)])}")),
                frames=len(trace) + 3), "exp03_hill_climbing_queens.gif", fps=1.5)
display(Image(filename=gif))
\end{lstlisting}
\noindent{\small\bfseries\color{gold!70!black}Output}
\begin{lstlisting}[style=out]
Solution: [6, 0, 2, 7, 5, 3, 1, 4] found after 1 random restarts;  attacking pairs = 0
\end{lstlisting}
\begin{lstlisting}[style=py]
hs = [attacks(t) for t in trace]
plt.figure(figsize=(6, 3.2)); plt.plot(hs, "-o", color=NAVY, mfc=GOLD, ms=8)
plt.xlabel("step"); plt.ylabel("attacking pairs h"); plt.title("Hill climbing descends h to 0 (successful run)")
plt.tight_layout(); plt.savefig(f"{FIG}/exp03_hc_curve.png", dpi=130); plt.show()
\end{lstlisting}
\begin{lstlisting}[style=py]
roads = {("Arad","Zerind"):75,("Arad","Sibiu"):140,("Arad","Timisoara"):118,("Zerind","Oradea"):71,("Oradea","Sibiu"):151,
 ("Timisoara","Lugoj"):111,("Lugoj","Mehadia"):70,("Mehadia","Drobeta"):75,("Drobeta","Craiova"):120,("Craiova","Rimnicu"):146,
 ("Craiova","Pitesti"):138,("Sibiu","Fagaras"):99,("Sibiu","Rimnicu"):80,("Rimnicu","Pitesti"):97,("Fagaras","Bucharest"):211,
 ("Pitesti","Bucharest"):101,("Bucharest","Giurgiu"):90,("Bucharest","Urziceni"):85}
G = {}
for (a, b), d in roads.items(): G.setdefault(a, {})[b] = d; G.setdefault(b, {})[a] = d
SLD = {"Arad":366,"Bucharest":0,"Craiova":160,"Drobeta":242,"Fagaras":176,"Giurgiu":77,"Lugoj":244,"Mehadia":241,
       "Oradea":380,"Pitesti":100,"Rimnicu":193,"Sibiu":253,"Timisoara":329,"Urziceni":80,"Zerind":374}
def best_first(start, goal, f):
    import heapq
    fr = [(f(start, 0), start, 0, [start])]; seen = set(); order = []
    while fr:
        _, n, g, path = heapq.heappop(fr)
        if n in seen: continue
        seen.add(n); order.append(n)
        if n == goal: return path, g, order
        for m, d in G[n].items(): heapq.heappush(fr, (f(m, g + d), m, g + d, path + [m]))
gp, gc, go = best_first("Arad", "Bucharest", lambda n, g: SLD[n])
ap, ac, ao = best_first("Arad", "Bucharest", lambda n, g: g + SLD[n])
print("Greedy best-first:", " -> ".join(gp), f"= {gc} km   (expanded {len(go)})")
print("A*               :", " -> ".join(ap), f"= {ac} km   (expanded {len(ao)})")
\end{lstlisting}
\noindent{\small\bfseries\color{gold!70!black}Output}
\begin{lstlisting}[style=out]
Greedy best-first: Arad -> Sibiu -> Fagaras -> Bucharest = 450 km   (expanded 4)
A*               : Arad -> Sibiu -> Rimnicu -> Pitesti -> Bucharest = 418 km   (expanded 6)
\end{lstlisting}
\subsection*{Output Figures and Animation}
\begin{center}
\animategraphics[controls,loop,autoplay,width=0.62\linewidth]{2}{frames/exp03_hill_climbing_queens/f-}{0}{3}\\
{\small\color{slate}\textbf{Animation} (4 frames; plays in Adobe Acrobat Reader). GIF: \texttt{assets/animations/exp03\_hill\_climbing\_queens.gif}}
\end{center}
\begin{center}
\includegraphics[width=0.78\linewidth,height=0.32\textheight,keepaspectratio]{../assets/figures/exp03_hc_curve.png}
\end{center}
\subsection*{Result}
\begin{outcome}[Result \& Conclusion]Plain steepest-ascent hill climbing solves only about one in seven 8-queens instances because it halts at local minima of h; random restarts always reach a solution. Greedy best-first reaches Bucharest quickly (450 km) but A* finds the optimal 418 km route.\end{outcome}
\subsection*{Viva-Voce Questions}
\begin{enumerate}
\item Why does hill climbing get stuck? Name three landscape features responsible.
\item How does random-restart hill climbing achieve completeness?
\item Why is greedy best-first search not optimal? Give the Romania counter-example.
\item What is the difference between hill climbing and simulated annealing?
\item What objective function did you use for 8-queens, and what is its maximum value?
\end{enumerate}
\newpage
\section*{Experiment 4: Game Playing --- Minimax with Alpha--Beta Pruning}\label{exp4}
\noindent\co{CO204.2}\ \btl{L3 Apply}\ \btl{L4 Analyse}\hfill\textbf{\color{navy}Date:}\ \rule{3cm}{0.4pt}\quad\textbf{\color{navy}Sign:}\ \rule{3cm}{0.4pt}\par\medskip
\begin{keypoint}[Aim]Implement an unbeatable Tic-Tac-Toe player using Minimax and Alpha--Beta pruning, and measure how many nodes pruning saves.\end{keypoint}
\subsection*{Theory}
In a two-player zero-sum game, \textbf{Minimax} backs up utilities: MAX takes the maximum and MIN the minimum of child values.

\textbf{Alpha--Beta} keeps $\alpha$ (best value for MAX) and $\beta$ (best value for MIN) and prunes a subtree as soon as $\alpha\ge\beta$. It returns exactly the same decision as Minimax; with perfect ordering, time drops from $O(b^m)$ to $O(b^{m/2})$.
\subsection*{Algorithm}
\begin{algo}[Procedure]\begin{enumerate}
\item If the state is terminal, return its utility (+1 X wins, --1 O wins, 0 draw).
\item MAX node: for each move, v = max(v, MIN-VALUE(child, $\alpha$, $\beta$)); if v $\ge$ $\beta$ prune; $\alpha$ = max($\alpha$, v).
\item MIN node: for each move, v = min(v, MAX-VALUE(child, $\alpha$, $\beta$)); if v $\le$ $\alpha$ prune; $\beta$ = min($\beta$, v).
\item At the root choose the move with the best backed-up value.
\end{enumerate}\end{algo}
\subsection*{Program (Python 3)}
{\small\color{slate}Notebook: \texttt{06\_Notebooks/Exp04\_Minimax\_AlphaBeta\_TicTacToe.ipynb}\\ Script: \texttt{06\_Notebooks/python/Exp04\_Minimax\_AlphaBeta\_TicTacToe.py}}
\begin{lstlisting}[style=py]
import math
LINES = [(0,1,2),(3,4,5),(6,7,8),(0,3,6),(1,4,7),(2,5,8),(0,4,8),(2,4,6)]
def winner(b):
    for i, j, k in LINES:
        if b[i] != " " and b[i] == b[j] == b[k]: return b[i]
    return "D" if " " not in b else None
def moves(b): return [i for i, v in enumerate(b) if v == " "]
def play(b, i, p): t = list(b); t[i] = p; return "".join(t)
def utility(w): return {"X": 1, "O": -1, "D": 0}[w]

counter = {"n": 0}
def minimax(b, player):
    counter["n"] += 1; w = winner(b)
    if w: return utility(w), None
    best = (-math.inf, None) if player == "X" else (math.inf, None)
    for m in moves(b):
        v, _ = minimax(play(b, m, player), "O" if player == "X" else "X")
        if (player == "X" and v > best[0]) or (player == "O" and v < best[0]): best = (v, m)
    return best

def alphabeta(b, player, a=-math.inf, be=math.inf):
    counter["n"] += 1; w = winner(b)
    if w: return utility(w), None
    best_m = None
    if player == "X":
        v = -math.inf
        for m in moves(b):
            val, _ = alphabeta(play(b, m, "X"), "O", a, be)
            if val > v: v, best_m = val, m
            a = max(a, v)
            if a >= be: break
    else:
        v = math.inf
        for m in moves(b):
            val, _ = alphabeta(play(b, m, "O"), "X", a, be)
            if val < v: v, best_m = val, m
            be = min(be, v)
            if a >= be: break
    return v, best_m

empty = " " * 9
counter["n"] = 0; v1, m1 = minimax(empty, "X"); n1 = counter["n"]
counter["n"] = 0; v2, m2 = alphabeta(empty, "X"); n2 = counter["n"]
print(f"Minimax    : value {v1}, best first move {m1}, nodes = {n1:,}")
print(f"Alpha-Beta : value {v2}, best first move {m2}, nodes = {n2:,}")
print(f"Pruning saved {100*(1-n2/n1):.1f}% of the nodes -- same decision, value 0 means perfect play is a draw.")
\end{lstlisting}
\noindent{\small\bfseries\color{gold!70!black}Output}
\begin{lstlisting}[style=out]
Minimax    : value 0, best first move 0, nodes = 549,946
Alpha-Beta : value 0, best first move 0, nodes = 18,297
Pruning saved 96.7% of the nodes -- same decision, value 0 means perfect play is a draw.
\end{lstlisting}
\begin{lstlisting}[style=py]
fig, ax = plt.subplots(figsize=(6, 3.4))
ax.bar(["Minimax", "Alpha-Beta"], [n1, n2], color=[SLATE, GOLD], edgecolor=NAVY)
for i, v in enumerate([n1, n2]): ax.text(i, v, f"{v:,}", ha="center", va="bottom", color=NAVY, weight="bold")
ax.set_ylabel("nodes generated"); ax.set_title("Tic-Tac-Toe from the empty board")
plt.tight_layout(); plt.savefig(f"{FIG}/exp04_nodes.png", dpi=130); plt.show()
\end{lstlisting}
\begin{lstlisting}[style=py]
b, p, history = empty, "X", [empty]
while not winner(b):
    _, m = alphabeta(b, p); b = play(b, m, p); history.append(b); p = "O" if p == "X" else "X"
print("Final board:"); [print(" " + " | ".join(b[r*3:r*3+3])) for r in range(3)]
print("Result:", {"D": "Draw", "X": "X wins", "O": "O wins"}[winner(b)])

fig, ax = plt.subplots(figsize=(4, 4.3))
def drawb(k):
    s = history[min(k, len(history) - 1)]; ax.clear(); ax.set_xlim(0, 3); ax.set_ylim(0, 3); ax.set_aspect("equal"); ax.axis("off")
    ax.add_patch(plt.Rectangle((0, 0), 3, 3, color=NAVY))
    for t in (1, 2): ax.plot([t, t], [0, 3], color=GOLD, lw=3); ax.plot([0, 3], [t, t], color=GOLD, lw=3)
    for i, v in enumerate(s):
        if v != " ": ax.text(i % 3 + .5, 2.5 - i // 3, v, ha="center", va="center", fontsize=40, weight="bold", color="white" if v == "X" else GOLD)
    ax.set_title(f"Move {min(k, len(history)-1)}")
gif = save_gif(animation.FuncAnimation(fig, drawb, frames=len(history) + 2), "exp04_tictactoe_selfplay.gif", fps=1)
display(Image(filename=gif))
\end{lstlisting}
\noindent{\small\bfseries\color{gold!70!black}Output}
\begin{lstlisting}[style=out]
Final board:
 X | X | O
 O | O | X
 X | O | X
Result: Draw
\end{lstlisting}
\subsection*{Output Figures and Animation}
\begin{center}
\animategraphics[controls,loop,autoplay,width=0.62\linewidth]{2}{frames/exp04_tictactoe_selfplay/f-}{0}{9}\\
{\small\color{slate}\textbf{Animation} (10 frames; plays in Adobe Acrobat Reader). GIF: \texttt{assets/animations/exp04\_tictactoe\_selfplay.gif}}
\end{center}
\begin{center}
\includegraphics[width=0.78\linewidth,height=0.32\textheight,keepaspectratio]{../assets/figures/exp04_nodes.png}
\end{center}
\subsection*{Result}
\begin{outcome}[Result \& Conclusion]Minimax and Alpha--Beta choose the same optimal move and agree that perfect play leads to a draw, but Alpha--Beta generates only a small fraction of the nodes. Self-play between two optimal players ends in a draw.\end{outcome}
\subsection*{Viva-Voce Questions}
\begin{enumerate}
\item What are $\alpha$ and $\beta$ and when is a branch pruned?
\item Does alpha--beta pruning ever change the Minimax decision? Why?
\item How does move ordering affect the effectiveness of pruning?
\item Why is full Minimax infeasible for chess?
\item What is an evaluation function and when is it used?
\end{enumerate}
\newpage
\section*{Experiment 5: Constraint Satisfaction --- N-Queens, Map Colouring and Sudoku}\label{exp5}
\noindent\co{CO204.2}\ \btl{L3 Apply}\ \btl{L4 Analyse}\hfill\textbf{\color{navy}Date:}\ \rule{3cm}{0.4pt}\quad\textbf{\color{navy}Sign:}\ \rule{3cm}{0.4pt}\par\medskip
\begin{keypoint}[Aim]Solve CSPs (N-Queens, map colouring of Australia and Sudoku) using backtracking search enhanced with MRV, forward checking and AC-3.\end{keypoint}
\subsection*{Theory}
A CSP is $\langle X,D,C\rangle$. \textbf{Backtracking} assigns one variable at a time and undoes an assignment when a variable has no legal value.

Heuristics: \textbf{MRV} (choose the variable with fewest legal values), \textbf{LCV}, \textbf{forward checking}. \textbf{AC-3} enforces arc consistency by deleting unsupported values.
\subsection*{Algorithm}
\begin{algo}[Procedure]\begin{enumerate}
\item If the assignment is complete, return it.
\item Select an unassigned variable (MRV).
\item For each value consistent with the assignment: assign it, run forward checking/AC-3, recurse.
\item If recursion fails, undo the assignment and try the next value; if none remain, backtrack.
\end{enumerate}\end{algo}
\subsection*{Program (Python 3)}
{\small\color{slate}Notebook: \texttt{06\_Notebooks/Exp05\_CSP\_NQueens\_MapColour\_Sudoku.ipynb}\\ Script: \texttt{06\_Notebooks/python/Exp05\_CSP\_NQueens\_MapColour\_Sudoku.py}}
\begin{lstlisting}[style=py]
def nqueens(n):
    sol, trace, calls = [], [], [0]
    def ok(col, row):
        return all(r != row and abs(r - row) != col - c for c, r in enumerate(sol))
    def bt(col):
        calls[0] += 1
        if col == n: return True
        for row in range(n):
            if ok(col, row):
                sol.append(row); trace.append(sol[:])
                if bt(col + 1): return True
                sol.pop(); trace.append(sol[:])
        return False
    bt(0); return sol, trace, calls[0]
sol, trace, calls = nqueens(8)
print("8-Queens solution (row of queen in each column):", sol, "| recursive calls:", calls, "| board states:", len(trace))
for n in (4, 6, 8, 10, 12):
    s, t, c = nqueens(n); print(f"  n={n:>2}: first solution {s}  calls={c}")
\end{lstlisting}
\noindent{\small\bfseries\color{gold!70!black}Output}
\begin{lstlisting}[style=out]
8-Queens solution (row of queen in each column): [0, 4, 7, 5, 2, 6, 1, 3] | recursive calls: 114 | board 
  n= 4: first solution [1, 3, 0, 2]  calls=9
  n= 6: first solution [1, 3, 5, 0, 2, 4]  calls=32
  n= 8: first solution [0, 4, 7, 5, 2, 6, 1, 3]  calls=114
  n=10: first solution [0, 2, 5, 7, 9, 4, 8, 1, 3, 6]  calls=103
  n=12: first solution [0, 2, 4, 7, 9, 11, 5, 10, 1, 6, 8, 3]  calls=262
\end{lstlisting}
\begin{lstlisting}[style=py]
fig, ax = plt.subplots(figsize=(4.5, 4.8))
keep = trace[::max(1, len(trace)//40)] + [trace[-1]] * 4
def drawq(k):
    s = keep[k]; ax.clear(); n = 8
    ax.imshow([[(i + j) % 2 for j in range(n)] for i in range(n)], cmap=matplotlib.colors.ListedColormap([MIST, "#C8D0DC"]))
    for c, r in enumerate(s): ax.text(c, r, "Q", ha="center", va="center", fontsize=24, fontfamily="DejaVu Sans", color=LEAF if len(s) == 8 else NAVY)
    ax.set_xticks([]); ax.set_yticks([]); ax.set_title(f"Backtracking: {len(s)} queens placed")
gif = save_gif(animation.FuncAnimation(fig, drawq, frames=len(keep)), "exp05_nqueens_backtracking.gif", fps=4)
display(Image(filename=gif))
\end{lstlisting}
\begin{lstlisting}[style=py]
NB = {"WA":["NT","SA"],"NT":["WA","SA","Q"],"SA":["WA","NT","Q","NSW","V"],"Q":["NT","SA","NSW"],"NSW":["Q","SA","V"],"V":["SA","NSW"],"T":[]}
COLS = ["red", "green", "blue"]
def colour(assign, domains, log):
    if len(assign) == len(NB): return assign
    var = min((v for v in NB if v not in assign), key=lambda v: (len(domains[v]), -len(NB[v])))   # MRV + degree
    for val in domains[var]:
        new_d = {k: list(d) for k, d in domains.items()}; new_d[var] = [val]; ok = True
        for nb in NB[var]:
            if nb not in assign and val in new_d[nb]:
                new_d[nb].remove(val); ok &= bool(new_d[nb])   # forward checking
        if ok:
            log.append(f"{var} = {val}"); r = colour({**assign, var: val}, new_d, log)
            if r: return r
            log.append(f"backtrack on {var}")
    return None
log = []; res = colour({}, {v: COLS[:] for v in NB}, log)
print("Assignment order:", " -> ".join(log)); print("Solution:", res)
\end{lstlisting}
\noindent{\small\bfseries\color{gold!70!black}Output}
\begin{lstlisting}[style=out]
Assignment order: SA = red -> NT = green -> Q = blue -> NSW = green -> WA = blue -> V = blue -> T = red
Solution: {'SA': 'red', 'NT': 'green', 'Q': 'blue', 'NSW': 'green', 'WA': 'blue', 'V': 'blue', 'T': 'red'
\end{lstlisting}
\begin{lstlisting}[style=py]
import networkx as nx
pos = {"WA":(0,1.2),"NT":(1.3,2),"SA":(1.6,0.9),"Q":(2.8,2),"NSW":(3,0.8),"V":(2.5,-0.2),"T":(2.7,-1.2)}
Gm = nx.Graph([(a, b) for a in NB for b in NB[a]]); Gm.add_node("T")
pal = {"red": CRIMSON, "green": LEAF, "blue": NAVY2}
plt.figure(figsize=(5.5, 4.2))
nx.draw(Gm, pos, with_labels=True, node_color=[pal[res[n]] for n in Gm.nodes], node_size=1700, font_color="white",
        font_weight="bold", font_family="serif", edge_color=GOLD, width=2.5)
plt.title("Australia map colouring -- constraint graph"); plt.savefig(f"{FIG}/exp05_map_colouring.png", dpi=130, bbox_inches="tight"); plt.show()
\end{lstlisting}
\begin{lstlisting}[style=py]
PUZ = "530070000600195000098000060800060003400803001700020006060000280000419005000080079"
cells = [(r, c) for r in range(9) for c in range(9)]
peers = {x: {y for y in cells if y != x and (y[0] == x[0] or y[1] == x[1] or (y[0]//3, y[1]//3) == (x[0]//3, x[1]//3))} for x in cells}
def ac3(D):
    q = [(x, y) for x in cells for y in peers[x]]
    while q:
        x, y = q.pop()
        if len(D[y]) == 1 and D[y][0] in D[x]:
            D[x] = [v for v in D[x] if v != D[y][0]]
            if not D[x]: return False
            q.extend((z, x) for z in peers[x] if z != y)
    return True
def solve(D, stats):
    if not ac3(D): return None
    if all(len(d) == 1 for d in D.values()): return D
    x = min((c for c in cells if len(D[c]) > 1), key=lambda c: len(D[c]))
    for v in D[x]:
        stats[0] += 1; E = {k: list(d) for k, d in D.items()}; E[x] = [v]
        r = solve(E, stats)
        if r: return r
    return None
D0 = {(r, c): [int(PUZ[r*9+c])] if PUZ[r*9+c] != "0" else list(range(1, 10)) for r, c in cells}
D1 = {k: list(v) for k, v in D0.items()}; ac3(D1)
print("Unknown cells before AC-3:", sum(len(v) > 1 for v in D0.values()), "| after AC-3 alone:", sum(len(v) > 1 for v in D1.values()))
stats = [0]; S = solve(D0, stats); print("Guesses needed by backtracking:", stats[0])
fig, ax = plt.subplots(figsize=(4.6, 4.6)); ax.set_xlim(0, 9); ax.set_ylim(0, 9); ax.axis("off")
for (r, c), v in S.items():
    given = PUZ[r*9+c] != "0"
    ax.add_patch(plt.Rectangle((c, 8 - r), 1, 1, color=GOLDL if given else "white", ec="#C8D0DC"))
    ax.text(c + .5, 8.5 - r, v[0], ha="center", va="center", fontsize=15, color=NAVY if given else TEAL, weight="bold" if given else "normal")
for t in range(0, 10, 3): ax.plot([t, t], [0, 9], color=NAVY, lw=2.5); ax.plot([0, 9], [t, t], color=NAVY, lw=2.5)
ax.set_title("Sudoku solved (givens shaded)")
HARD = "800000000003600000070090200050007000000045700000100030001000068008500010090000400"   # 'world's hardest' (Inkala)
Dh = {(r, c): [int(HARD[r*9+c])] if HARD[r*9+c] != "0" else list(range(1, 10)) for r, c in cells}
Dh1 = {k: list(v) for k, v in Dh.items()}; ac3(Dh1)
st2 = [0]; Sh = solve(Dh, st2)
print(f"Hard puzzle: unknown after AC-3 = {sum(len(v) > 1 for v in Dh1.values())}, backtracking guesses = {st2[0]}")
print("Hard puzzle first row solved:", [Sh[(0, c)][0] for c in range(9)]); plt.savefig(f"{FIG}/exp05_sudoku.png", dpi=130, bbox_inches="tight"); plt.show()
\end{lstlisting}
\noindent{\small\bfseries\color{gold!70!black}Output}
\begin{lstlisting}[style=out]
Unknown cells before AC-3: 51 | after AC-3 alone: 0
Guesses needed by backtracking: 0
Hard puzzle: unknown after AC-3 = 60, backtracking guesses = 1850
Hard puzzle first row solved: [8, 1, 2, 7, 5, 3, 6, 4, 9]
\end{lstlisting}
\subsection*{Output Figures and Animation}
\begin{center}
\animategraphics[controls,loop,autoplay,width=0.62\linewidth]{2}{frames/exp05_nqueens_backtracking/f-}{0}{14}\\
{\small\color{slate}\textbf{Animation} (15 frames; plays in Adobe Acrobat Reader). GIF: \texttt{assets/animations/exp05\_nqueens\_backtracking.gif}}
\end{center}
\begin{center}
\includegraphics[width=0.78\linewidth,height=0.32\textheight,keepaspectratio]{../assets/figures/exp05_map_colouring.png}
\end{center}
\begin{center}
\includegraphics[width=0.78\linewidth,height=0.32\textheight,keepaspectratio]{../assets/figures/exp05_sudoku.png}
\end{center}
\subsection*{Result}
\begin{outcome}[Result \& Conclusion]Backtracking finds the first 8-queens solution in about a hundred recursive calls; MRV with forward checking colours Australia without any backtracking; AC-3 alone solves an easy Sudoku, while the 'hardest' Sudoku additionally needs backtracking guesses guided by MRV.\end{outcome}
\subsection*{Viva-Voce Questions}
\begin{enumerate}
\item Define a CSP with its three components for map colouring.
\item What is the MRV heuristic and why is it called fail-first?
\item Differentiate forward checking from arc consistency.
\item What is the complexity of AC-3?
\item Why is Tasmania irrelevant in the Australia map-colouring problem?
\end{enumerate}
\newpage
\section*{Experiment 6: Facts, Rules and Logical Inference}\label{exp6}
\noindent\co{CO204.3}\ \btl{L3 Apply}\ \btl{L4 Analyse}\hfill\textbf{\color{navy}Date:}\ \rule{3cm}{0.4pt}\quad\textbf{\color{navy}Sign:}\ \rule{3cm}{0.4pt}\par\medskip
\begin{keypoint}[Aim]Represent facts and rules in a knowledge base (Prolog-style, in Python) and answer queries using forward chaining and backward chaining, including a propositional truth-table entailment checker.\end{keypoint}
\subsection*{Theory}
A \textbf{knowledge base} stores facts (ground atoms) and rules (Horn clauses $p_1\wedge\dots\wedge p_n\Rightarrow q$).

\textbf{Forward chaining} fires rules whose premises are satisfied until no new facts appear (data-driven).

\textbf{Backward chaining} starts from the query and recursively proves the rule premises (goal-driven, as in Prolog).

\textbf{Model checking} decides $KB\models\alpha$ by enumerating all truth assignments.
\subsection*{Algorithm}
\begin{algo}[Procedure]\begin{enumerate}
\item Store facts as tuples, e.g. ('parent', 'ram', 'sita'); rules as (head, [body literals]) with variables starting with '?'.
\item Forward chaining: repeatedly unify rule bodies with known facts and add instantiated heads.
\item Backward chaining: to prove a goal, unify it with a fact or a rule head and prove the body recursively.
\item Print the derived facts, query answers and proof traces.
\end{enumerate}\end{algo}
\subsection*{Program (Python 3)}
{\small\color{slate}Notebook: \texttt{06\_Notebooks/Exp06\_Logic\_Facts\_Rules\_Inference.ipynb}\\ Script: \texttt{06\_Notebooks/python/Exp06\_Logic\_Facts\_Rules\_Inference.py}}
\begin{lstlisting}[style=py]
import itertools
FACTS = {("parent","dashrath","ram"), ("parent","dashrath","bharat"), ("parent","ram","luv"), ("parent","ram","kush"),
         ("parent","kaushalya","ram"), ("male","dashrath"), ("male","ram"), ("male","bharat"), ("male","luv"), ("male","kush"),
         ("female","kaushalya"), ("female","sita"), ("parent","sita","luv"), ("parent","sita","kush")}
RULES = [
  (("father","?x","?y"),      [("parent","?x","?y"), ("male","?x")]),
  (("mother","?x","?y"),      [("parent","?x","?y"), ("female","?x")]),
  (("grandparent","?x","?z"), [("parent","?x","?y"), ("parent","?y","?z")]),
  (("sibling","?x","?y"),     [("parent","?p","?x"), ("parent","?p","?y"), ("neq","?x","?y")]),
  (("uncle","?u","?c"),       [("parent","?p","?c"), ("sibling","?u","?p"), ("male","?u")]),
]
def is_var(t): return isinstance(t, str) and t.startswith("?")
def walk(t, th):
    while is_var(t) and t in th: t = th[t]
    return t
def subst(lit, th): return tuple(walk(t, th) for t in lit)
def match(pat, fact, th):
    if len(pat) != len(fact): return None
    th = dict(th)
    for p, f in zip(pat, fact):
        p = th.get(p, p)
        if is_var(p): th[p] = f
        elif p != f: return None
    return th
print(f"KB: {len(FACTS)} facts, {len(RULES)} rules")
\end{lstlisting}
\noindent{\small\bfseries\color{gold!70!black}Output}
\begin{lstlisting}[style=out]
KB: 14 facts, 5 rules
\end{lstlisting}
\begin{lstlisting}[style=py]
def forward_chain(facts, rules, verbose=True):
    facts = set(facts); rnd = 0; history = []
    while True:
        rnd += 1; new = set()
        for head, body in rules:
            def sat(i, th):
                if i == len(body): yield th; return
                lit = subst(body[i], th)
                if lit[0] == "neq":
                    if lit[1] != lit[2]: yield from sat(i + 1, th)
                    return
                for f in facts:
                    t2 = match(lit, f, th)
                    if t2 is not None: yield from sat(i + 1, t2)
            for th in sat(0, {}):
                h = subst(head, th)
                if h not in facts: new.add(h)
        if not new: break
        history.append((rnd, sorted(new))); facts |= new
        if verbose: print(f"Round {rnd}: +{len(new)} facts ->", ", ".join(f"{p}({', '.join(a)})" for p, *a in sorted(new)))
    return facts, history
ALL, hist = forward_chain(FACTS, RULES)
print("\nGrandparents of luv:", sorted(f[1] for f in ALL if f[0] == "grandparent" and f[2] == "luv"))
print("Uncle of luv       :", sorted(f[1] for f in ALL if f[0] == "uncle" and f[2] == "luv"))
\end{lstlisting}
\noindent{\small\bfseries\color{gold!70!black}Output}
\begin{lstlisting}[style=out]
Round 1: +15 facts -> father(dashrath, bharat), father(dashrath, ram), father(ram, kush), father(ram, luv
Round 2: +2 facts -> uncle(bharat, kush), uncle(bharat, luv)

Grandparents of luv: ['dashrath', 'kaushalya']
Uncle of luv       : ['bharat']
\end{lstlisting}
\begin{lstlisting}[style=py]
counter = itertools.count()
def rename(rule):
    k = next(counter); ren = lambda l: tuple(f"{t}_{k}" if is_var(t) else t for t in l)
    return ren(rule[0]), [ren(b) for b in rule[1]]
def bc(goals, th, depth=0, trace=None):
    if not goals: yield th; return
    g = subst(goals[0], th); rest = goals[1:]
    if g[0] == "neq":
        if g[1] != g[2]: yield from bc(rest, th, depth, trace)
        return
    for f in sorted(FACTS):
        t2 = match(g, f, th)
        if t2 is not None:
            if trace is not None: trace.append("  " * depth + f"[ok] fact {f[0]}({', '.join(f[1:])})")
            yield from bc(rest, t2, depth, trace)
    for rule in RULES:
        head, body = rename(rule)
        # unify the renamed rule head with the goal
        t2 = dict(th); ok = True
        for a, b in zip(head, g):
            a, b = walk(a, t2), walk(b, t2)
            if a == b: continue
            if is_var(a): t2[a] = b
            elif is_var(b): t2[b] = a
            else: ok = False; break
        if ok and len(head) == len(g):
            if trace is not None: trace.append("  " * depth + f"-> try rule {rule[0][0]} for {g[0]}({', '.join(g[1:])})")
            yield from bc(body + rest, t2, depth + 1, trace)
def ask(q):
    tr = []; sols = []
    for th in bc([q], {}, 0, tr):
        sols.append(subst(q, th)); break
    return sols, tr
for q in [("grandparent","dashrath","kush"), ("uncle","bharat","luv"), ("mother","sita","kush"), ("father","bharat","luv")]:
    s, tr = ask(q)
    print(f"?- {q[0]}({', '.join(q[1:])}).  ->  {'YES' if s else 'NO'}")
s, tr = ask(("uncle","bharat","luv")); print("\nProof trace for uncle(bharat, luv):"); print("\n".join(tr[:14]))
\end{lstlisting}
\noindent{\small\bfseries\color{gold!70!black}Output}
\begin{lstlisting}[style=out]
?- grandparent(dashrath, kush).  ->  YES
?- uncle(bharat, luv).  ->  YES
?- mother(sita, kush).  ->  YES
?- father(bharat, luv).  ->  NO

Proof trace for uncle(bharat, luv):
-> try rule uncle for uncle(bharat, luv)
  [ok] fact parent(ram, luv)
  -> try rule sibling for sibling(bharat, ram)
    [ok] fact parent(dashrath, bharat)
    [ok] fact parent(dashrath, ram)
    [ok] fact male(bharat)
\end{lstlisting}
\begin{lstlisting}[style=py]
import pandas as pd
syms = ["P", "Q", "R"]
KB    = lambda m: (not m["P"] or m["Q"]) and (not m["Q"] or m["R"]) and m["P"]   # P=>Q, Q=>R, P
alpha = lambda m: m["R"]
rows = []
for vals in itertools.product([False, True], repeat=3):
    m = dict(zip(syms, vals)); rows.append({**{k: "T" if v else "F" for k, v in m.items()}, "KB": "T" if KB(m) else "F", "alpha": "T" if alpha(m) else "F"})
df = pd.DataFrame(rows); df.columns = ["P", "Q", "R", "KB", "alpha = R"]; print(df.to_string(index=False))
entails = all(alpha(dict(zip(syms, v))) for v in itertools.product([False, True], repeat=3) if KB(dict(zip(syms, v))))
print("\nKB |= R ?", entails)
\end{lstlisting}
\noindent{\small\bfseries\color{gold!70!black}Output}
\begin{lstlisting}[style=out]
P Q R KB alpha = R
F F F  F     F
F F T  F     T
F T F  F     F
F T T  F     T
T F F  F     F
T F T  F     T
T T F  F     F
T T T  T     T

KB |= R ? True
\end{lstlisting}
\begin{lstlisting}[style=py]
import networkx as nx
G = nx.DiGraph(); base = [f for f in FACTS if f[0] == "parent"]
for f in base: G.add_edge(f[1], f[2], kind="parent")
for f in ALL:
    if f[0] in ("grandparent", "uncle"): G.add_edge(f[1], f[2], kind=f[0])
pos = {"dashrath": (0.5, 2), "kaushalya": (2, 2), "ram": (1.2, 1), "bharat": (-0.6, 1), "sita": (2.8, 1), "luv": (0.6, 0), "kush": (2.2, 0)}
fig, ax = plt.subplots(figsize=(7, 5))
def drawkb(k):
    ax.clear(); kinds = ["parent"] + (["grandparent"] if k >= 1 else []) + (["uncle"] if k >= 2 else [])
    cmap = {"parent": SLATE, "grandparent": GOLD, "uncle": CRIMSON}
    es = [(u, v) for u, v, d in G.edges(data=True) if d["kind"] in kinds]
    nx.draw_networkx_nodes(G, pos, ax=ax, node_color=NAVY, node_size=1400)
    nx.draw_networkx_labels(G, pos, ax=ax, font_color="white", font_family="serif", font_size=9)
    nx.draw_networkx_edges(G, pos, ax=ax, edgelist=es, edge_color=[cmap[G[u][v]["kind"]] for u, v in es], width=2.2, arrowsize=14,
                           connectionstyle="arc3,rad=0.15", node_size=1400)
    ax.set_xlim(-1.2, 3.4); ax.set_ylim(-0.5, 2.5)
    ax.set_title(["Base facts: parent", "Forward chaining adds grandparent (gold)", "...then uncle (crimson)"][min(k, 2)]); ax.axis("off")
gif = save_gif(animation.FuncAnimation(fig, drawkb, frames=4), "exp06_forward_chaining.gif", fps=0.8)
display(Image(filename=gif))
\end{lstlisting}
\subsection*{Output Figures and Animation}
\begin{center}
\animategraphics[controls,loop,autoplay,width=0.62\linewidth]{2}{frames/exp06_forward_chaining/f-}{0}{2}\\
{\small\color{slate}\textbf{Animation} (3 frames; plays in Adobe Acrobat Reader). GIF: \texttt{assets/animations/exp06\_forward\_chaining.gif}}
\end{center}
\subsection*{Result}
\begin{outcome}[Result \& Conclusion]Forward chaining derived all father, mother, grandparent, sibling and uncle relations in a few rounds; backward chaining answered individual queries with an explicit proof trace; model checking confirmed that {P$\Rightarrow$Q, Q$\Rightarrow$R, P} entails R.\end{outcome}
\subsection*{Viva-Voce Questions}
\begin{enumerate}
\item Differentiate forward and backward chaining.
\item What is a Horn clause and why is it important for chaining?
\item How does Prolog implement backward chaining?
\item What is the complexity of truth-table model checking?
\item Why did we need the 'neq' built-in in the sibling rule?
\end{enumerate}
\newpage
\section*{Experiment 7: Unification and Resolution in First-Order Logic}\label{exp7}
\noindent\co{CO204.3}\ \btl{L3 Apply}\ \btl{L4 Analyse}\hfill\textbf{\color{navy}Date:}\ \rule{3cm}{0.4pt}\quad\textbf{\color{navy}Sign:}\ \rule{3cm}{0.4pt}\par\medskip
\begin{keypoint}[Aim]Implement the unification algorithm (with occurs check) and a resolution refutation prover for first-order clauses; prove classic textbook theorems.\end{keypoint}
\subsection*{Theory}
\textbf{Unification} finds the most general substitution $\theta$ such that $p\theta=q\theta$. The \textbf{occurs check} rejects bindings like $x/f(x)$.

\textbf{Resolution:} from clauses $(\ell_1\vee\alpha)$ and $(\neg\ell_2\vee\beta)$ with $\textsc{Unify}(\ell_1,\ell_2)=\theta$ derive $(\alpha\vee\beta)\theta$.

To prove $KB\models q$: add $\neg q$ to the clause set and derive the empty clause $\square$.
\subsection*{Algorithm}
\begin{algo}[Procedure]\begin{enumerate}
\item Represent terms as nested tuples: variables are lowercase strings starting with '?'; f(x) is ('f', x).
\item UNIFY(x, y, $\theta$): if equal return $\theta$; if x is a variable bind it (with occurs check); if compound unify arguments pairwise.
\item Convert KB and negated goal into clauses (sets of literals).
\item Repeatedly resolve pairs of clauses; stop when the empty clause appears (proof) or no new clauses (no proof).
\end{enumerate}\end{algo}
\subsection*{Program (Python 3)}
{\small\color{slate}Notebook: \texttt{06\_Notebooks/Exp07\_Unification\_Resolution\_FOL.ipynb}\\ Script: \texttt{06\_Notebooks/python/Exp07\_Unification\_Resolution\_FOL.py}}
\begin{lstlisting}[style=py]
def is_var(t): return isinstance(t, str) and t.startswith("?")
def walk(t, th):
    while is_var(t) and t in th: t = th[t]
    return t
def occurs(v, t, th):
    t = walk(t, th)
    if t == v: return True
    return isinstance(t, tuple) and any(occurs(v, a, th) for a in t[1:])
def unify(x, y, th=None):
    if th is None: th = {}
    x, y = walk(x, th), walk(y, th)
    if x == y: return th
    if is_var(x): return None if occurs(x, y, th) else {**th, x: y}
    if is_var(y): return unify(y, x, th)
    if isinstance(x, tuple) and isinstance(y, tuple) and x[0] == y[0] and len(x) == len(y):
        for a, b in zip(x[1:], y[1:]):
            th = unify(a, b, th)
            if th is None: return None
        return th
    return None
def resolve_term(t, th):
    t = walk(t, th)
    return (t[0],) + tuple(resolve_term(a, th) for a in t[1:]) if isinstance(t, tuple) else t
def show(t):
    if isinstance(t, tuple): return f"{t[0]}({', '.join(show(a) for a in t[1:])})"
    return t.lstrip("?")
def show_th(th): return "{" + ", ".join(f"{k.lstrip('?')}/{show(resolve_term(v, th))}" for k, v in th.items()) + "}" if th is not None else "FAIL"

tests = [(("Knows","John","?x"), ("Knows","John","Jane")),
         (("Knows","John","?x"), ("Knows","?y","Bill")),
         (("Knows","John","?x"), ("Knows","?y",("Mother","?y"))),
         (("Knows","John","?x"), ("Knows","?x","Elizabeth")),
         (("P","?x",("f","?y")), ("P","a",("f",("g","?z")))),
         (("Q","?x"), ("Q",("f","?x")))]
for a, b in tests: print(f"UNIFY({show(a):<22}, {show(b):<26}) = {show_th(unify(a, b))}")
\end{lstlisting}
\noindent{\small\bfseries\color{gold!70!black}Output}
\begin{lstlisting}[style=out]
UNIFY(Knows(John, x)        , Knows(John, Jane)         ) = {x/Jane}
UNIFY(Knows(John, x)        , Knows(y, Bill)            ) = {y/John, x/Bill}
UNIFY(Knows(John, x)        , Knows(y, Mother(y))       ) = {y/John, x/Mother(John)}
UNIFY(Knows(John, x)        , Knows(x, Elizabeth)       ) = FAIL
UNIFY(P(x, f(y))            , P(a, f(g(z)))             ) = {x/a, y/g(z)}
UNIFY(Q(x)                  , Q(f(x))                   ) = FAIL
\end{lstlisting}
\begin{lstlisting}[style=py]
import itertools
def neg(l): return l[1] if l[0] == "not" else ("not", l)
def subst_lit(l, th): return ("not", resolve_term(l[1], th)) if l[0] == "not" else resolve_term(l, th)
cid = itertools.count()
def standardise(clause):
    k = next(cid); m = {}
    def rn(t):
        if is_var(t): m.setdefault(t, f"{t}{k}"); return m[t]
        return (t[0],) + tuple(rn(a) for a in t[1:]) if isinstance(t, tuple) else t
    return frozenset(("not", rn(l[1])) if l[0] == "not" else rn(l) for l in clause)
def resolvents(c1, c2):
    out = []
    for l1 in c1:
        for l2 in c2:
            if (l1[0] == "not") != (l2[0] == "not"):
                a, b = (l1[1], l2) if l1[0] == "not" else (l1, l2[1])
                th = unify(a, b)
                if th is not None:
                    new = frozenset(subst_lit(l, th) for l in (c1 - {l1}) | (c2 - {l2}))
                    out.append((new, th, l1, l2))
    return out
def prove(kb, goal, max_iter=2000):
    clauses = [standardise(c) for c in kb] + [standardise([neg(goal)])]
    names = {c: f"C{i+1}" for i, c in enumerate(clauses)}; steps = []; seen = set(clauses)
    for _ in range(max_iter):
        new = []
        for c1, c2 in itertools.combinations(list(seen), 2):
            for r, th, l1, l2 in resolvents(c1, c2):
                if r in seen or r in [n for n, *_ in new]: continue
                new.append((r, c1, c2, th))
                if not r:
                    steps.append((r, c1, c2, th)); return True, steps, names
        if not new: return False, steps, names
        for r, c1, c2, th in sorted(new, key=lambda x: len(x[0]))[:30]:
            seen.add(r); names[r] = f"C{len(names)+1}"; steps.append((r, c1, c2, th))
    return False, steps, names
def cs(c): return " | ".join(("~" + show(l[1])) if l[0] == "not" else show(l) for l in sorted(c, key=str)) or "[]"

# Example 1 -- Socrates
KB1 = [[("not", ("Man","?x")), ("Mortal","?x")], [("Man","Socrates")]]
ok, st, nm = prove(KB1, ("Mortal","Socrates"))
print("KB: forall x Man(x) => Mortal(x);  Man(Socrates)   Goal: Mortal(Socrates)")
for r, c1, c2, th in st: print(f"   resolve [{cs(c1)}] with [{cs(c2)}]  =>  {cs(r)}")
print("Proved!" if ok else "Not proved")
\end{lstlisting}
\noindent{\small\bfseries\color{gold!70!black}Output}
\begin{lstlisting}[style=out]
KB: forall x Man(x) => Mortal(x);  Man(Socrates)   Goal: Mortal(Socrates)
   resolve [Mortal(x0) | ~Man(x0)] with [Man(Socrates)]  =>  Mortal(Socrates)
   resolve [Mortal(x0) | ~Man(x0)] with [~Mortal(Socrates)]  =>  ~Man(Socrates)
   resolve [~Mortal(Socrates)] with [Mortal(Socrates)]  =>  []
Proved!
\end{lstlisting}
\begin{lstlisting}[style=py]
# Example 2 -- Colonel West (Russell & Norvig 9.3)
KB2 = [
 [("not",("American","?x")),("not",("Weapon","?y")),("not",("Sells","?x","?y","?z")),("not",("Hostile","?z")),("Criminal","?x")],
 [("not",("Missile","?x")),("not",("Owns","Nono","?x")),("Sells","West","?x","Nono")],
 [("not",("Missile","?x")),("Weapon","?x")],
 [("not",("Enemy","?x","America")),("Hostile","?x")],
 [("Owns","Nono","M1")],[("Missile","M1")],[("American","West")],[("Enemy","Nono","America")]]
ok, st, nm = prove(KB2, ("Criminal","West"))
print("Goal: Criminal(West) ->", "PROVED [ok]" if ok else "not proved", f"after generating {len(st)} resolvents")
chain = []
r = st[-1]
print("Final step:", f"[{cs(r[1])}] + [{cs(r[2])}] => {cs(r[0])}")
\end{lstlisting}
\noindent{\small\bfseries\color{gold!70!black}Output}
\begin{lstlisting}[style=out]
Goal: Criminal(West) -> PROVED [ok] after generating 100 resolvents
Final step: [~Sells(West, M1, Nono)] + [Sells(West, M1, Nono)] => []
\end{lstlisting}
\begin{lstlisting}[style=py]
# Visualise the Socrates refutation as a proof tree
fig, ax = plt.subplots(figsize=(7.5, 4)); ax.axis("off")
boxes = {"C1": (0.15, 0.85, "~Man(x) | Mortal(x)"), "C3": (0.62, 0.85, "~Mortal(Socrates)  [negated goal]"),
         "R1": (0.38, 0.5, "~Man(Socrates)   theta={x/Socrates}"), "C2": (0.85, 0.5, "Man(Socrates)"), "R2": (0.6, 0.12, "[]  (empty clause)")}
for k, (x, y, t) in boxes.items():
    ax.text(x, y, t, ha="center", va="center", fontsize=12, color="white" if k.startswith("R") else NAVY,
            bbox=dict(boxstyle="round,pad=0.5", fc=NAVY if k.startswith("R") else GOLDL, ec=GOLD, lw=2))
for a, b in (("C1","R1"),("C3","R1"),("R1","R2"),("C2","R2")):
    ax.annotate("", xy=boxes[b][:2], xytext=boxes[a][:2], arrowprops=dict(arrowstyle="-|>", color=TEAL, lw=2, shrinkA=18, shrinkB=18))
ax.set_title("Resolution refutation proof of Mortal(Socrates)")
plt.savefig(f"{FIG}/exp07_resolution_tree.png", dpi=130, bbox_inches="tight"); plt.show()
\end{lstlisting}
\subsection*{Output Figures and Animation}
\begin{center}
\includegraphics[width=0.78\linewidth,height=0.32\textheight,keepaspectratio]{../assets/figures/exp07_resolution_tree.png}
\end{center}
\subsection*{Result}
\begin{outcome}[Result \& Conclusion]The unifier returned the expected MGUs, correctly failed on the clash requiring standardising apart, and the occurs check blocked x/f(x). The resolution prover derived the empty clause for both Mortal(Socrates) and Criminal(West).\end{outcome}
\subsection*{Viva-Voce Questions}
\begin{enumerate}
\item What is a most general unifier?
\item Why is the occurs check needed?
\item Why must variables be standardised apart before resolution?
\item What is Skolemisation?
\item Why is resolution called refutation-complete?
\end{enumerate}
\newpage
\section*{Experiment 8: Bayes' Theorem and a Bayesian Network for Diagnosis}\label{exp8}
\noindent\co{CO204.4}\ \btl{L3 Apply}\ \btl{L4 Analyse}\hfill\textbf{\color{navy}Date:}\ \rule{3cm}{0.4pt}\quad\textbf{\color{navy}Sign:}\ \rule{3cm}{0.4pt}\par\medskip
\begin{keypoint}[Aim]Apply Bayes' theorem to a medical-test diagnosis problem and implement the Burglary--Alarm Bayesian network with exact inference by enumeration and approximate inference by sampling.\end{keypoint}
\subsection*{Theory}
\textbf{Bayes' theorem:} $P(H\mid E)=\dfrac{P(E\mid H)P(H)}{P(E)}$.

A \textbf{Bayesian network} is a DAG with a CPT per node; the joint is $\prod_i P(x_i\mid parents(X_i))$.

\textbf{Enumeration:} $P(X\mid e)=\alpha\sum_y P(X,e,y)$. \textbf{Rejection sampling} estimates the same posterior from random samples.
\subsection*{Algorithm}
\begin{algo}[Procedure]\begin{enumerate}
\item Encode priors and conditional probabilities.
\item Compute the posterior with Bayes' rule and study how it varies with the prior.
\item Encode the Burglary network CPTs; compute the joint probability of a full assignment.
\item Compute P(Burglary | JohnCalls, MaryCalls) by enumeration and compare with rejection sampling.
\end{enumerate}\end{algo}
\subsection*{Program (Python 3)}
{\small\color{slate}Notebook: \texttt{06\_Notebooks/Exp08\_Bayes\_BayesianNetwork\_Diagnosis.ipynb}\\ Script: \texttt{06\_Notebooks/python/Exp08\_Bayes\_BayesianNetwork\_Diagnosis.py}}
\begin{lstlisting}[style=py]
def bayes(prior, sens, fpr):
    evidence = sens * prior + fpr * (1 - prior)
    return sens * prior / evidence, evidence
post, ev = bayes(prior=0.01, sens=0.90, fpr=0.05)
print(f"P(+) = {ev:.4f}")
print(f"P(Disease | +) = {post:.4f}  ->  only {post:.1%} despite a 90%-sensitive test (base-rate effect)")
post2, _ = bayes(prior=post, sens=0.90, fpr=0.05)
print(f"After a second independent positive test: {post2:.1%}")
\end{lstlisting}
\noindent{\small\bfseries\color{gold!70!black}Output}
\begin{lstlisting}[style=out]
P(+) = 0.0585
P(Disease | +) = 0.1538  ->  only 15.4% despite a 90%-sensitive test (base-rate effect)
After a second independent positive test: 76.6%
\end{lstlisting}
\begin{lstlisting}[style=py]
priors = np.linspace(0.001, 0.5, 300)
plt.figure(figsize=(6.8, 3.8))
for sens, fpr, c in ((0.90, 0.05, NAVY), (0.99, 0.01, TEAL), (0.80, 0.10, CRIMSON)):
    plt.plot(priors, [bayes(p, sens, fpr)[0] for p in priors], color=c, lw=2.5, label=f"sens={sens}, FPR={fpr}")
plt.axvline(0.01, color=GOLD, ls="--"); plt.text(0.015, 0.05, "prevalence 1%", color=GOLD)
plt.xlabel("prior P(Disease)"); plt.ylabel("posterior P(Disease | +)"); plt.legend(frameon=False); plt.title("Posterior depends strongly on the prior")
plt.tight_layout(); plt.savefig(f"{FIG}/exp08_posterior_vs_prior.png", dpi=130); plt.show()
\end{lstlisting}
\begin{lstlisting}[style=py]
import itertools
P_B, P_E = 0.001, 0.002
P_A = {(True, True): 0.95, (True, False): 0.94, (False, True): 0.29, (False, False): 0.001}
P_J = {True: 0.90, False: 0.05}; P_M = {True: 0.70, False: 0.01}
def p(v, prob): return prob if v else 1 - prob
def joint(b, e, a, j, m):
    return p(b, P_B) * p(e, P_E) * p(a, P_A[(b, e)]) * p(j, P_J[a]) * p(m, P_M[a])
print(f"P(j, m, a, ~b, ~e) = {joint(False, False, True, True, True):.6f}")

def enumerate_query(var, evidence):
    names = ["B", "E", "A", "J", "M"]; dist = {}
    for val in (True, False):
        tot = 0.0
        for vals in itertools.product((True, False), repeat=5):
            asg = dict(zip(names, vals))
            if asg[var] != val or any(asg[k] != v for k, v in evidence.items()): continue
            tot += joint(*vals)
        dist[val] = tot
    z = sum(dist.values()); return {k: v / z for k, v in dist.items()}
for q, ev in (("B", {"J": True, "M": True}), ("E", {"J": True, "M": True}), ("A", {"J": True, "M": True}), ("B", {"J": True, "M": False})):
    r = enumerate_query(q, ev); print(f"P({q}=T | {', '.join(f'{k}={v}' for k, v in ev.items())}) = {r[True]:.4f}")
\end{lstlisting}
\noindent{\small\bfseries\color{gold!70!black}Output}
\begin{lstlisting}[style=out]
P(j, m, a, ~b, ~e) = 0.000628
P(B=T | J=True, M=True) = 0.2842
P(E=T | J=True, M=True) = 0.1761
P(A=T | J=True, M=True) = 0.7607
P(B=T | J=True, M=False) = 0.0051
\end{lstlisting}
\begin{lstlisting}[style=py]
rng = np.random.default_rng(0)
def prior_sample():
    b = rng.random() < P_B; e = rng.random() < P_E; a = rng.random() < P_A[(b, e)]
    return b, e, a, rng.random() < P_J[a], rng.random() < P_M[a]
def likelihood_weighting(n):
    wsum = {True: 0.0, False: 0.0}; est = []
    for i in range(1, n + 1):
        b = rng.random() < P_B; e = rng.random() < P_E; a = rng.random() < P_A[(b, e)]
        w = P_J[a] * P_M[a]; wsum[b] += w
        if i % 500 == 0: est.append(wsum[True] / (wsum[True] + wsum[False]))
    return est
exact = enumerate_query("B", {"J": True, "M": True})[True]
est = likelihood_weighting(100_000)
print(f"Exact P(B | j, m) = {exact:.4f};  likelihood weighting (100k samples) ~= {est[-1]:.4f}")
fig, ax = plt.subplots(figsize=(6.8, 3.6)); xs = np.arange(1, len(est) + 1) * 500
line, = ax.plot([], [], color=NAVY, lw=2); ax.axhline(exact, color=GOLD, ls="--", lw=2, label=f"exact = {exact:.3f}")
ax.set_xlim(0, xs[-1]); ax.set_ylim(0, 0.8); ax.set_xlabel("samples"); ax.set_ylabel("estimate"); ax.legend(frameon=False)
def upd(k):
    line.set_data(xs[:k + 1], est[:k + 1]); ax.set_title(f"Likelihood weighting converges -- {xs[k]:,} samples"); return line,
gif = save_gif(animation.FuncAnimation(fig, upd, frames=range(0, len(est), 5)), "exp08_sampling_convergence.gif", fps=6)
display(Image(filename=gif))
\end{lstlisting}
\noindent{\small\bfseries\color{gold!70!black}Output}
\begin{lstlisting}[style=out]
Exact P(B | j, m) = 0.2842;  likelihood weighting (100k samples) ~= 0.2975
\end{lstlisting}
\begin{lstlisting}[style=py]
import networkx as nx
G = nx.DiGraph([("Burglary","Alarm"),("Earthquake","Alarm"),("Alarm","JohnCalls"),("Alarm","MaryCalls")])
pos = {"Burglary":(0,1),"Earthquake":(2,1),"Alarm":(1,0.5),"JohnCalls":(0,0),"MaryCalls":(2,0)}
plt.figure(figsize=(6, 3.6))
nx.draw(G, pos, with_labels=True, node_color=[GOLDL, GOLDL, NAVY, MIST, MIST], node_size=3800, edgecolors=NAVY,
        font_color="black", font_family="serif", font_size=10, arrowsize=18, edge_color=TEAL, width=2)
plt.title("Bayesian network structure"); plt.savefig(f"{FIG}/exp08_bayesnet.png", dpi=130, bbox_inches="tight"); plt.show()
\end{lstlisting}
\subsection*{Output Figures and Animation}
\begin{center}
\animategraphics[controls,loop,autoplay,width=0.62\linewidth]{2}{frames/exp08_sampling_convergence/f-}{0}{14}\\
{\small\color{slate}\textbf{Animation} (15 frames; plays in Adobe Acrobat Reader). GIF: \texttt{assets/animations/exp08\_sampling\_convergence.gif}}
\end{center}
\begin{center}
\includegraphics[width=0.78\linewidth,height=0.32\textheight,keepaspectratio]{../assets/figures/exp08_bayesnet.png}
\end{center}
\begin{center}
\includegraphics[width=0.78\linewidth,height=0.32\textheight,keepaspectratio]{../assets/figures/exp08_posterior_vs_prior.png}
\end{center}
\subsection*{Result}
\begin{outcome}[Result \& Conclusion]A positive result from a 90\%-sensitive test with 1\% prevalence gives only a 15.4\% chance of disease. In the alarm network, P(Burglary | JohnCalls, MaryCalls) = 0.284 by exact enumeration, and likelihood weighting converges to the same value.\end{outcome}
\subsection*{Viva-Voce Questions}
\begin{enumerate}
\item State Bayes' theorem and name each term.
\item Explain the base-rate fallacy with the medical test example.
\item Why is a Bayesian network more compact than the full joint distribution?
\item What is the complexity of inference by enumeration?
\item How does likelihood weighting differ from rejection sampling?
\end{enumerate}
\newpage
\section*{Experiment 9: Decision-Tree Classifier and Rule Extraction}\label{exp9}
\noindent\co{CO204.5}\ \btl{L3 Apply}\ \btl{L4 Analyse}\hfill\textbf{\color{navy}Date:}\ \rule{3cm}{0.4pt}\quad\textbf{\color{navy}Sign:}\ \rule{3cm}{0.4pt}\par\medskip
\begin{keypoint}[Aim]Build a decision tree (ID3 from scratch on the Play-Tennis data and scikit-learn CART on the Iris data), interpret the learned rules and study overfitting with depth.\end{keypoint}
\subsection*{Theory}
Entropy $H(S)=-\sum_k p_k\log_2p_k$; information gain $Gain(S,A)=H(S)-\sum_v\frac{|S_v|}{|S|}H(S_v)$.

ID3 chooses the attribute with the highest gain at each node and recurses until nodes are pure. CART uses the Gini index. Deep trees overfit; limit depth or prune.
\subsection*{Algorithm}
\begin{algo}[Procedure]\begin{enumerate}
\item Compute the entropy of the full dataset.
\item Compute the information gain of every attribute; split on the best.
\item Recurse on each subset until pure or no attributes remain.
\item Convert each root-to-leaf path into an IF--THEN rule; evaluate on test data.
\end{enumerate}\end{algo}
\subsection*{Program (Python 3)}
{\small\color{slate}Notebook: \texttt{06\_Notebooks/Exp09\_DecisionTree\_Rules.ipynb}\\ Script: \texttt{06\_Notebooks/python/Exp09\_DecisionTree\_Rules.py}}
\begin{lstlisting}[style=py]
import pandas as pd
from math import log2
data = pd.DataFrame([
 ["Sunny","Hot","High","Weak","No"],["Sunny","Hot","High","Strong","No"],["Overcast","Hot","High","Weak","Yes"],
 ["Rain","Mild","High","Weak","Yes"],["Rain","Cool","Normal","Weak","Yes"],["Rain","Cool","Normal","Strong","No"],
 ["Overcast","Cool","Normal","Strong","Yes"],["Sunny","Mild","High","Weak","No"],["Sunny","Cool","Normal","Weak","Yes"],
 ["Rain","Mild","Normal","Weak","Yes"],["Sunny","Mild","Normal","Strong","Yes"],["Overcast","Mild","High","Strong","Yes"],
 ["Overcast","Hot","Normal","Weak","Yes"],["Rain","Mild","High","Strong","No"]],
 columns=["Outlook","Temperature","Humidity","Wind","Play"])
def H(s):
    p = s.value_counts(normalize=True); return -sum(q * log2(q) for q in p if q > 0)
def gain(df, a, t="Play"):
    return H(df[t]) - sum(len(g) / len(df) * H(g[t]) for _, g in df.groupby(a))
print(f"H(S) = {H(data['Play']):.3f}")
gains = {a: gain(data, a) for a in data.columns[:-1]}
for a, g in sorted(gains.items(), key=lambda x: -x[1]): print(f"  Gain(S, {a:<11}) = {g:.3f}")
\end{lstlisting}
\noindent{\small\bfseries\color{gold!70!black}Output}
\begin{lstlisting}[style=out]
H(S) = 0.940
  Gain(S, Outlook    ) = 0.247
  Gain(S, Humidity   ) = 0.152
  Gain(S, Wind       ) = 0.048
  Gain(S, Temperature) = 0.029
\end{lstlisting}
\begin{lstlisting}[style=py]
def id3(df, attrs, t="Play"):
    if df[t].nunique() == 1: return df[t].iloc[0]
    if not attrs: return df[t].mode()[0]
    best = max(attrs, key=lambda a: gain(df, a))
    return {best: {v: id3(g, [a for a in attrs if a != best]) for v, g in df.groupby(best)}}
tree = id3(data, list(data.columns[:-1]))
def rules(node, cond=()):
    if not isinstance(node, dict): yield cond, node; return
    (a, branches), = node.items()
    for v, sub in branches.items(): yield from rules(sub, cond + ((a, v),))
print("Learned tree:", tree, "\n")
for c, out in rules(tree): print("IF " + " AND ".join(f"{a} = {v}" for a, v in c) + f"  THEN Play = {out}")
def predict(node, x):
    while isinstance(node, dict): (a, br), = node.items(); node = br[x[a]]
    return node
acc = np.mean([predict(tree, r) == r["Play"] for _, r in data.iterrows()]); print(f"\nTraining accuracy: {acc:.0%}")
\end{lstlisting}
\noindent{\small\bfseries\color{gold!70!black}Output}
\begin{lstlisting}[style=out]
Learned tree: {'Outlook': {'Overcast': 'Yes', 'Rain': {'Wind': {'Strong': 'No', 'Weak': 'Yes'}}, 'Sunny':

IF Outlook = Overcast  THEN Play = Yes
IF Outlook = Rain AND Wind = Strong  THEN Play = No
IF Outlook = Rain AND Wind = Weak  THEN Play = Yes
IF Outlook = Sunny AND Humidity = High  THEN Play = No
IF Outlook = Sunny AND Humidity = Normal  THEN Play = Yes

Training accuracy: 100%
\end{lstlisting}
\begin{lstlisting}[style=py]
fig, ax = plt.subplots(figsize=(6.5, 3.4))
ks = sorted(gains, key=lambda k: -gains[k])
ax.bar(ks, [gains[k] for k in ks], color=[NAVY, TEAL, GOLD, CRIMSON])
for i, k in enumerate(ks): ax.text(i, gains[k] + 0.005, f"{gains[k]:.3f}", ha="center", color=NAVY, weight="bold")
ax.set_ylabel("information gain (bits)"); ax.set_title("Root selection: Outlook has the highest gain")
plt.tight_layout(); plt.savefig(f"{FIG}/exp09_gain.png", dpi=130); plt.show()
\end{lstlisting}
\begin{lstlisting}[style=py]
from sklearn.datasets import load_iris
from sklearn.model_selection import train_test_split
from sklearn.tree import DecisionTreeClassifier, export_text, plot_tree
from sklearn.metrics import accuracy_score, confusion_matrix
iris = load_iris(); X, y = iris.data, iris.target
Xtr, Xte, ytr, yte = train_test_split(X, y, test_size=0.3, random_state=42, stratify=y)
clf = DecisionTreeClassifier(max_depth=3, random_state=0).fit(Xtr, ytr)
print("Test accuracy:", round(accuracy_score(yte, clf.predict(Xte)), 3))
print("Confusion matrix:\n", confusion_matrix(yte, clf.predict(Xte)))
print("\nRules learned:\n" + export_text(clf, feature_names=list(iris.feature_names)))
plt.figure(figsize=(10, 5.5))
plot_tree(clf, feature_names=iris.feature_names, class_names=iris.target_names, filled=True, rounded=True, fontsize=9)
plt.title("CART decision tree (max_depth = 3)"); plt.savefig(f"{FIG}/exp09_iris_tree.png", dpi=130, bbox_inches="tight"); plt.show()
\end{lstlisting}
\noindent{\small\bfseries\color{gold!70!black}Output}
\begin{lstlisting}[style=out]
Test accuracy: 0.933
Confusion matrix:
 [[15  0  0]
 [ 0 14  1]
 [ 0  2 13]]

Rules learned:
|--- petal width (cm) <= 0.70
|   |--- class: 0
|--- petal width (cm) >  0.70
|   |--- petal width (cm) <= 1.55
|   |   |--- petal length (cm) <= 4.95
|   |   |   |--- class: 1
|   |   |--- petal length (cm) >  4.95
|   |   |   |--- class: 2
|   |--- petal width (cm) >  1.55
|   |   |--- petal length (cm) <= 4.85
|   |   |   |--- class: 1
|   |   |--- petal length (cm) >  4.85
|   |   |   |--- class: 2
\end{lstlisting}
\begin{lstlisting}[style=py]
depths = range(1, 11); tr_acc, te_acc = [], []
for d in depths:
    m = DecisionTreeClassifier(max_depth=d, random_state=0).fit(Xtr, ytr)
    tr_acc.append(m.score(Xtr, ytr)); te_acc.append(m.score(Xte, yte))
plt.figure(figsize=(6.5, 3.5)); plt.plot(depths, tr_acc, "-o", color=NAVY, label="train"); plt.plot(depths, te_acc, "-s", color=GOLD, label="test")
plt.xlabel("max depth"); plt.ylabel("accuracy"); plt.legend(frameon=False); plt.title("Depth vs accuracy -- watch for overfitting")
plt.tight_layout(); plt.savefig(f"{FIG}/exp09_depth.png", dpi=130); plt.show()

X2 = X[:, 2:4]; xx, yy = np.meshgrid(np.linspace(0.5, 7.2, 300), np.linspace(0, 2.7, 300))
fig, ax = plt.subplots(figsize=(6, 4.4)); cm = matplotlib.colors.ListedColormap([GOLDL, "#BFE3E6", "#D9DEEA"])
def frm(k):
    d = k + 1; m = DecisionTreeClassifier(max_depth=d, random_state=0).fit(X2, y); ax.clear()
    ax.contourf(xx, yy, m.predict(np.c_[xx.ravel(), yy.ravel()]).reshape(xx.shape), cmap=cm)
    for c, col in zip(range(3), [GOLD, TEAL, NAVY]): ax.scatter(*X2[y == c].T, color=col, edgecolor="white", s=40, label=iris.target_names[c])
    ax.set_xlabel("petal length"); ax.set_ylabel("petal width"); ax.legend(frameon=False, loc="upper left")
    ax.set_title(f"Decision boundary, depth = {d}, train acc = {m.score(X2, y):.2f}")
gif = save_gif(animation.FuncAnimation(fig, frm, frames=6), "exp09_tree_boundary.gif", fps=1)
display(Image(filename=gif))
\end{lstlisting}
\subsection*{Output Figures and Animation}
\begin{center}
\animategraphics[controls,loop,autoplay,width=0.62\linewidth]{2}{frames/exp09_tree_boundary/f-}{0}{5}\\
{\small\color{slate}\textbf{Animation} (6 frames; plays in Adobe Acrobat Reader). GIF: \texttt{assets/animations/exp09\_tree\_boundary.gif}}
\end{center}
\begin{center}
\includegraphics[width=0.78\linewidth,height=0.32\textheight,keepaspectratio]{../assets/figures/exp09_depth.png}
\end{center}
\begin{center}
\includegraphics[width=0.78\linewidth,height=0.32\textheight,keepaspectratio]{../assets/figures/exp09_gain.png}
\end{center}
\begin{center}
\includegraphics[width=0.78\linewidth,height=0.32\textheight,keepaspectratio]{../assets/figures/exp09_iris_tree.png}
\end{center}
\subsection*{Result}
\begin{outcome}[Result \& Conclusion]ID3 selected Outlook (gain 0.247) as the root and produced five readable rules with 100\% training accuracy. CART on Iris reached high test accuracy at depth 3; deeper trees fit training data perfectly without improving test accuracy --- a sign of overfitting.\end{outcome}
\subsection*{Viva-Voce Questions}
\begin{enumerate}
\item Define entropy and information gain.
\item Why is Outlook chosen as the root for Play Tennis?
\item Compare Gini index and entropy.
\item How does pruning reduce overfitting?
\item Why are decision trees considered interpretable models?
\end{enumerate}
\newpage
\section*{Experiment 10: Rule-Based Expert System and a Basic Chatbot (Mini-Project)}\label{exp10}
\noindent\co{CO204.5}\ \btl{L3 Apply}\ \btl{L4 Analyse}\hfill\textbf{\color{navy}Date:}\ \rule{3cm}{0.4pt}\quad\textbf{\color{navy}Sign:}\ \rule{3cm}{0.4pt}\par\medskip
\begin{keypoint}[Aim]Develop a rule-based expert system with certainty factors and an explanation facility for crop-disease advice, and a pattern-matching chatbot for college admission queries.\end{keypoint}
\subsection*{Theory}
An expert system has a \textbf{knowledge base} (IF--THEN rules), \textbf{working memory} (facts), an \textbf{inference engine} (match $\to$ conflict resolution $\to$ act) and an \textbf{explanation facility}.

Certainty factors combine as $CF(conclusion)=\min(CF_{premises})\times CF_{rule}$ and, for two rules supporting the same conclusion, $CF_1+CF_2(1-CF_1)$.

A \textbf{retrieval chatbot} matches user input to intents using patterns/keywords and returns a response.
\subsection*{Algorithm}
\begin{algo}[Procedure]\begin{enumerate}
\item Load rules and the user's observed facts with certainty values.
\item Find all rules whose premises are in working memory (match).
\item Fire them, computing the conclusion CF; combine CFs for repeated conclusions; repeat until no change.
\item Report conclusions sorted by CF and answer WHY/HOW questions from the firing trace.
\end{enumerate}\end{algo}
\subsection*{Program (Python 3)}
{\small\color{slate}Notebook: \texttt{06\_Notebooks/Exp10\_ExpertSystem\_Chatbot.ipynb}\\ Script: \texttt{06\_Notebooks/python/Exp10\_ExpertSystem\_Chatbot.py}}
\begin{lstlisting}[style=py]
RULES = [
 ("R1", ["yellow_leaves", "leaf_spots"],          "fungal_infection",   0.7),
 ("R2", ["white_powder_on_leaves"],                "powdery_mildew",     0.9),
 ("R3", ["fungal_infection", "humid_weather"],     "leaf_blight",        0.8),
 ("R4", ["wilting", "dry_soil"],                   "water_stress",       0.85),
 ("R5", ["wilting", "brown_stem_base"],            "root_rot",           0.75),
 ("R6", ["holes_in_leaves", "larvae_seen"],        "pest_attack",        0.9),
 ("R7", ["leaf_blight"],                           "advise_fungicide",   0.9),
 ("R8", ["powdery_mildew"],                        "advise_sulphur_spray", 0.85),
 ("R9", ["water_stress"],                          "advise_irrigation",  0.95),
 ("R10",["pest_attack"],                           "advise_neem_oil",    0.8),
 ("R11",["root_rot"],                              "advise_drainage",    0.8),
 ("R12",["yellow_leaves", "stunted_growth"],       "nitrogen_deficiency",0.7),
 ("R13",["nitrogen_deficiency"],                   "advise_urea_dose",   0.8),
]
def infer(observed):
    wm = dict(observed); trace = []; changed = True
    while changed:
        changed = False
        for rid, prem, concl, cf in RULES:
            if all(p in wm for p in prem):
                new = round(min(wm[p] for p in prem) * cf, 3)
                old = wm.get(concl)
                comb = new if old is None else round(old + new * (1 - old), 3)
                if (rid, concl) not in [(t[0], t[2]) for t in trace]:
                    wm[concl] = comb; trace.append((rid, prem, concl, comb)); changed = True
    return wm, trace
def explain(concl, trace):
    for rid, prem, c, cf in trace:
        if c == concl: return f"{concl} (CF {cf}) because {rid}: " + " AND ".join(prem)
    return f"{concl} was observed directly."
case = {"yellow_leaves": 1.0, "leaf_spots": 0.9, "humid_weather": 0.8, "holes_in_leaves": 0.6, "larvae_seen": 0.7}
wm, trace = infer(case)
print("Observed:", case, "\n\nInference trace:")
for rid, prem, c, cf in trace: print(f"  {rid}: IF {' AND '.join(prem)} THEN {c}   -> CF = {cf}")
print("\nAdvice (sorted by certainty):")
for k, v in sorted(wm.items(), key=lambda x: -x[1]):
    if k.startswith("advise"): print(f"  * {k.replace('advise_', '').replace('_', ' ').title():<16} CF = {v}")
print("\nHOW was leaf_blight concluded?  ", explain("leaf_blight", trace))
\end{lstlisting}
\noindent{\small\bfseries\color{gold!70!black}Output}
\begin{lstlisting}[style=out]
Observed: {'yellow_leaves': 1.0, 'leaf_spots': 0.9, 'humid_weather': 0.8, 'holes_in_leaves': 0.6, 'larvae

Inference trace:
  R1: IF yellow_leaves AND leaf_spots THEN fungal_infection   -> CF = 0.63
  R3: IF fungal_infection AND humid_weather THEN leaf_blight   -> CF = 0.504
  R6: IF holes_in_leaves AND larvae_seen THEN pest_attack   -> CF = 0.54
  R7: IF leaf_blight THEN advise_fungicide   -> CF = 0.454
  R10: IF pest_attack THEN advise_neem_oil   -> CF = 0.432

Advice (sorted by certainty):
  * Fungicide        CF = 0.454
  * Neem Oil         CF = 0.432

HOW was leaf_blight concluded?   leaf_blight (CF 0.504) because R3: fungal_infection AND humid_weather
\end{lstlisting}
\begin{lstlisting}[style=py]
import networkx as nx
G = nx.DiGraph()
for rid, prem, c, cf in trace:
    for p in prem: G.add_edge(p, c, label=rid)
layers = {}
def depth(n): return 0 if n in case else 1 + max(depth(p) for p in G.predecessors(n))
for n in G.nodes: layers.setdefault(depth(n), []).append(n)
pos = {n: (d * 3.0, -i * 1.2 + len(ns) * 0.6) for d, ns in layers.items() for i, n in enumerate(sorted(ns))}
order = [rid for rid, *_ in trace]
fig, ax = plt.subplots(figsize=(9, 4.8))
def fr(k):
    ax.clear(); fired = set(order[:k]); lit = set(case) | {c for rid, _, c, _ in trace if rid in fired}
    nx.draw_networkx_edges(G, pos, ax=ax, edge_color=[GOLD if G[u][v]["label"] in fired else "#C8D0DC" for u, v in G.edges], width=2.2, arrowsize=14)
    nx.draw_networkx_nodes(G, pos, ax=ax, node_color=[(LEAF if n.startswith("advise") else NAVY) if n in lit else MIST for n in G.nodes], node_size=2600, edgecolors=NAVY)
    for group, fc in ((lit, "white"), (set(G.nodes) - lit, NAVY)):
        nx.draw_networkx_labels(G, pos, ax=ax, labels={n: n.replace("_", "\n") for n in G.nodes if n in group}, font_size=7.5,
                                font_family="serif", font_color=fc)
    ax.set_title(f"Rule firing: {', '.join(order[:k]) or 'initial facts'}"); ax.axis("off"); ax.margins(0.12)
gif = save_gif(animation.FuncAnimation(fig, fr, frames=len(order) + 3), "exp10_expert_system_firing.gif", fps=1)
display(Image(filename=gif))
\end{lstlisting}
\begin{lstlisting}[style=py]
import re
INTENTS = [
 ("greet",    r"\b(hi|hello|hey|namaste)\b",                 "Hello! I am the Sanjivani AI&DS admission assistant. How can I help?"),
 ("fees",     r"\b(fee|fees|cost|tuition)\b",                "Fee details for B.Tech AI & DS are on the university admissions portal; scholarships are available for eligible students."),
 ("eligible", r"\b(eligib\w*|criteria|qualify|jee|mht.?cet)\b", "Eligibility: 10+2 with Physics & Mathematics and a valid MHT-CET / JEE Main score."),
 ("courses",  r"\b(course|subject|syllabus|curriculum)\b",   "The curriculum covers AI, ML, Deep Learning, Data Engineering and an industry internship."),
 ("hostel",   r"\b(hostel|accommodation|stay)\b",            "Separate hostels for boys and girls are available on campus."),
 ("placement",r"\b(placement\w*|jobs?|package|recruit\w*)\b",     "The Training & Placement cell runs aptitude training and campus drives every year."),
 ("contact",  r"\b(contact|phone|email|address)\b",          "Please visit the university website or the admissions office, Kopargaon."),
 ("bye",      r"\b(bye|thanks|thank you)\b",                 "Thank you! All the best for your admission."),
]
def reply(msg):
    m = msg.lower()
    for name, pat, ans in INTENTS:
        if re.search(pat, m): return name, ans
    return "fallback", "Sorry, I did not understand. You can ask about fees, eligibility, courses, hostel or placements."
for q in ["Hi there!", "What is the eligibility for AI and DS?", "How much are the fees?", "Is hostel available?",
          "Tell me about placements", "Can I study astrophysics?", "Thanks, bye"]:
    intent, ans = reply(q); print(f"User: {q}\nBot: [{intent}] {ans}\n")
\end{lstlisting}
\noindent{\small\bfseries\color{gold!70!black}Output}
\begin{lstlisting}[style=out]
User: Hi there!
Bot: [greet] Hello! I am the Sanjivani AI&DS admission assistant. How can I help?

User: What is the eligibility for AI and DS?
Bot: [eligible] Eligibility: 10+2 with Physics & Mathematics and a valid MHT-CET / JEE Main score.

User: How much are the fees?
Bot: [fees] Fee details for B.Tech AI & DS are on the university admissions portal; scholarships are avai

User: Is hostel available?
Bot: [hostel] Separate hostels for boys and girls are available on campus.

User: Tell me about placements
Bot: [placement] The Training & Placement cell runs aptitude training and campus drives every year.

User: Can I study astrophysics?
Bot: [fallback] Sorry, I did not understand. You can ask about fees, eligibility, courses, hostel or plac

User: Thanks, bye
Bot: [bye] Thank you! All the best for your admission.
\end{lstlisting}
\begin{lstlisting}[style=py]
tests = [("hello", "greet"), ("fees structure?", "fees"), ("am I eligible with JEE", "eligible"), ("syllabus please", "courses"),
         ("where do I stay", "hostel"), ("average package", "placement"), ("email id", "contact"), ("thank you", "bye"), ("weather today", "fallback")]
acc = np.mean([reply(q)[0] == t for q, t in tests]); print(f"Intent accuracy on test prompts: {acc:.0%}")
\end{lstlisting}
\noindent{\small\bfseries\color{gold!70!black}Output}
\begin{lstlisting}[style=out]
Intent accuracy on test prompts: 100%
\end{lstlisting}
\subsection*{Output Figures and Animation}
\begin{center}
\animategraphics[controls,loop,autoplay,width=0.62\linewidth]{2}{frames/exp10_expert_system_firing/f-}{0}{5}\\
{\small\color{slate}\textbf{Animation} (6 frames; plays in Adobe Acrobat Reader). GIF: \texttt{assets/animations/exp10\_expert\_system\_firing.gif}}
\end{center}
\subsection*{Result}
\begin{outcome}[Result \& Conclusion]The expert system chained symptom rules to diagnoses and advice, combined certainty factors and explained its reasoning; the chatbot correctly mapped all test prompts to intents, falling back gracefully for unknown queries.\end{outcome}
\subsection*{Viva-Voce Questions}
\begin{enumerate}
\item List the components of an expert system.
\item How are certainty factors combined in MYCIN?
\item What is conflict resolution in a production system?
\item What is the role of the explanation facility?
\item How would you improve the chatbot using machine learning?
\end{enumerate}
\end{document}